\chapter{Theory}

% ----------------------------------------------------------------------
\section{Force Constants}
\label{sec:fd}
% ----------------------------------------------------------------------

\subsection{Potential-energy expansion}

For theoretical description of lattice vibrations,  the crystal potential energy
in the atomic displacements is expanded from from equilibrium
\cite{born_dynamical_1996, maradudin_theory_1963},
\begin{multline}\label{eq:e_exp_ph}
    E\bigl(\{u_i\}_{i=1}^{\ndof}\bigr)
    = E_0 +
      \sum_{i}^{\ndof} \frac{\partial E}{\partial u_i}\, u_i +
      \frac{1}{2}\sum_{ij}^{\ndof}
        \frac{\partial^2 E}{\partial u_i\,\partial u_j}\, u_i\, u_j \\
    + \frac{1}{6}\sum_{ijk}^{\ndof}
        \frac{\partial^3 E}{\partial u_i\,\partial u_j\,\partial u_k}\,
        u_i\, u_j\, u_k
    + \cdots,
\end{multline}
where $\ndof = 3\nsuper = 3\Ncell\nprim$ is the total number of
degrees of freedom (see \cref{sec:notation_fourier} for the notation).
At equilibrium, the first-order term vanishes.  The
second-order coefficients define the harmonic, or second-order,
interatomic force constants (IFC2), the third-order coefficients
define the cubic, or third-order, force constants (IFC3).
The harmonic approximation retains
\cref{eq:e_exp_ph} to second order, and is sufficient for phonon
dispersion and most equilibrium thermodynamic quantities
\cite{dove_introduction_1993, wallace_thermodynamics_1998}.
Anharmonic lattice-dynamics methods include the cubic and potentially higher-order terms.
These introduce phonon--phonon scattering processes and hence the finite linewidth,
frequency shift, and thermal transport properties of crystalline solids
\cite{ziman_electrons_2001, srivastava_physics_2019}.

% ----------------------------------------------------------------------
\subsection{Harmonic force constants (IFC2)}
\label{sub:fc2}
% ----------------------------------------------------------------------

Introducing the primitive-cell degree-of-freedom (DOF) index
$\mu = (\kappa, \alpha)$, which bundles the basis atom $\kappa$ and
the Cartesian direction $\alpha$ into a single label, the harmonic
force constants are defined as
\begin{equation}\label{eq:fc2_def}
  \Phi(0\mu,\, l'\mu')
  = \frac{\partial^2 E}
         {\partial u_{0\mu}\,\partial u_{l'\mu'}},
\end{equation}
where $u_{l\mu}$ is the displacement of DOF $\mu$ in cell $l$.
Within a perfect crystal, translational symmetry lets us fix the
first cell index to the home cell ($l{=}0$) without loss of
generality, so that $\Phi(0\mu,\,l'\mu')$ depends only on the
relative cell separation.  Grouping the three Cartesian directions
for a fixed pair of atoms $(\kappa,\kappa')$ yields a
$3\times 3$ block whose rows and columns correspond to force and
displacement components, respectively.

The IFC2 quantify how a small displacement of DOF $\mu'$ in cell
$l'$ changes the force on DOF $\mu$ in the home cell.  Linearising
the Newtonian equations of motion around equilibrium then gives the
force--displacement relation
\begin{equation}\label{eq:force-disp}
  F(0\mu)
  = -\sum_{l'\mu'} \Phi(0\mu,\, l'\mu')\, u_{l'\mu'},
\end{equation}
which is crucial to both harmonic lattice dynamics and the
finite-displacement schemes used to compute the IFC2 discussed in
\cref{sub:fc_computation}.

By construction, the IFC2 tensor obeys three exact relations that are
useful both for reducing storage and for enforcing physical
consistency when force constants are obtained numerically.
The first is permutation symmetry: since partial derivatives
commute, the tensor is symmetric under exchange of its two legs,
\begin{equation}\label{eq:fc2_perm}
  \Phi(0\mu,\,l'\mu') = \Phi(l'\mu',\,0\mu),
\end{equation}
which, combined with translational invariance, can equivalently be
written as
\begin{equation}
  \Phi(0\mu,\,l'\mu')
  =
  \Phi(0\mu',\,(-l')\mu).
  \label{eq:ifc-symmetry-equivalent}
\end{equation}
This form is commonly used in lattice-dynamical treatments of force
constants
\cite{born_dynamical_1996, maradudin_theory_1963, gonze_dynamical_1997}.

The second symmetry is translational invariance: a rigid translation of
the whole crystal in any Cartesian direction $\alpha'$ cannot generate
a restoring force, which leads to the acoustic sum rule
\begin{equation}\label{eq:fc2_asr}
  \sum_{l'\kappa'} \Phi_{\alpha\alpha'}(0\kappa,\,l'\kappa') = 0
  \qquad \forall\;\kappa,\alpha,\alpha',
\end{equation}
where the condition must hold for each triple
$(\kappa,\alpha,\alpha')$ separately \cite{born_dynamical_1996, pick_microscopic_1970, gonze_dynamical_1997}.
This relation guarantees that the three
acoustic branches have zero frequency at the $\Gamma$
point, and in practice its enforcement is critical for numerically
computed IFC2, as even small residual violations produce spurious
acoustic gaps in the dispersion.  The third is space-group
symmetry: under each operation $\{S|\vc{v}\}$ of the crystal point
group, the IFC2 transform as a second-rank Cartesian tensor on the
index pair and simultaneously map one atom--cell pair onto another.
The resulting linear relations among the entries of $\Phi$
typically reduce the number of independent parameters by one or two
orders of magnitude, and are exploited systematically by software
packages such as phonopy \cite{togo_first_2015} and symfc
\cite{seko_projector-based_2024}, and are discussed in-depth in \citet{fu_group_2019}.

Phonon frequencies and polarisation vectors are obtained as
eigenpairs of the mass-weighted Fourier transform of $\Phi$, the
dynamical matrix (see \cref{sec:notation_fourier} for the two phase
conventions used in this work).  See the standard literature, e.g.
\citet{born_dynamical_1996}, \citet{maradudin_theory_1963}, and
\citet{dove_introduction_1993}.

% ----------------------------------------------------------------------
\subsection{Anharmonic force constants (IFC3)}
\label{sub:fc3}
% ----------------------------------------------------------------------

The cubic, or third-order, interatomic force constants are the
generalisation of \cref{eq:fc2_def} to the next order of the
Taylor expansion \cref{eq:e_exp_ph}.  In the composite-index
notation, they read
\begin{equation}\label{eq:fc3_def}
  \Phi^{(3)}(0\mu,\,l'\mu',\,l''\mu'')
  = \frac{\partial^3 E}
         {\partial u_{0\mu}\,
          \partial u_{l'\mu'}\,
          \partial u_{l''\mu''}},
\end{equation}
and, as before, translational symmetry lets us fix the first cell
index to the home cell.  The IFC3 tensor captures the leading
anharmonic corrections to lattice dynamics.  It sets the amplitude
of three-phonon scattering processes, determines the finite
linewidth and temperature-dependent frequency shift of phonon modes,
and provides the dominant resistive mechanism for lattice thermal
conductivity in most crystalline solids \cite{ziman_electrons_2001, srivastava_physics_2019}.

The symmetry properties of IFC3 parallel those of IFC2, but with a
richer permutation structure due to the three legs of the
tensor.  Commutativity of partial derivatives now gives full $S_3$
symmetry, so that $\Phi^{(3)}(0\mu,\,l'\mu',\,l''\mu'')$ is
invariant under all six simultaneous permutations of its three
cell--DOF pairs:
\begin{equation}\label{eq:fc3_perm}
  \Phi^{(3)}(0\mu,\,l'\mu',\,l''\mu'')
  = \Phi^{(3)}(l'\mu',\,0\mu,\,l''\mu'')
  = \Phi^{(3)}(l''\mu'',\,l'\mu',\,0\mu)
  = \cdots
\end{equation}
\cite{srivastava_physics_2019}.  Translational
invariance again generates acoustic sum rules, now one per leg:
summing $\Phi^{(3)}$ over the cell--atom labels of any one leg must
vanish for all choices of the remaining indices,
\begin{equation}\label{eq:fc3_asr}
  \sum_{l''\kappa''}
     \Phi^{(3)}_{\alpha\alpha'\alpha''}(0\kappa,\,l'\kappa',\,l''\kappa'')
  = 0
  \qquad \forall\;\kappa,\kappa',l',\alpha,\alpha',\alpha'',
\end{equation}
and analogously under summation over the first or second leg
\cite{srivastava_physics_2019}.  As for the
harmonic case, enforcing these sum rules on numerically extracted
IFC3 is essential for a consistent phonon self-energy and, hence,
for accurate thermal-transport predictions \cite{li_shengbte_2014}.
Space-group symmetry acts as a third-rank Cartesian tensor on the
index triple while simultaneously permuting the atom--cell labels,
and typically reduces the independent-entry count by one to two
orders of magnitude \cite{fu_group_2019, seko_projector-based_2024,
esfarjani_method_2008}.

Despite these reductions, the storage requirements for the full IFC3
grow quickly with system size.  After fixing one cell index by
translational symmetry, the number of independent entries scales as
$\cO(\nprim\,\ndof^2) = \cO(\Ncell^2\nprim^3)$, which rapidly
becomes prohibitive for large supercells.  This scaling,
together with the even more unfavourable
cost of the four-tensor contraction that enters the phonon--phonon
self-energy, motivates the compression schemes developed in
\cref{sub:fc3_compression}.

% ----------------------------------------------------------------------
\subsection{Computation of force constants}
\label{sub:fc_computation}
% ----------------------------------------------------------------------

Several established methods exist for obtaining interatomic force
constants from first-principles calculations.  We summarise the
main approaches below. The choice among them involves trade-offs
between computational cost, accessible interaction range, the
order of anharmonicity that can be treated, and availability of software.

% ----------------------------------------------------------------------
\subsubsection{Finite displacements}
\label{ssub:fd_fc}
% ----------------------------------------------------------------------
The most widely used approach constructs a supercell, displaces atoms
one at a time along symmetry-inequivalent directions, performs a
self-consistent DFT calculation for each displaced structure, and
reconstructs the force-constant tensor by inverting the
force--displacement relation \cref{eq:force-disp}.  The method was
introduced by \citet{parlinski_first-principles_1997} and is
implemented in codes such as phonopy \cite{togo_first_2015},
PHON \cite{alfe_phon_2009}, and Phonon
\cite{parlinski_first-principles_1997}.  Space-group symmetry reduces
the number of independent displacements and hence the number of
required DFT calculations, making the approach practical for cells of
moderate size.

For cubic force constants, the analogous procedure requires displacing
two atoms simultaneously. The number of required supercell
calculations therefore grows quadratically with the number of atoms in
the primitive cell
\cite{esfarjani_method_2008, li_shengbte_2014}.  \Citet{esfarjani_method_2008}
introduced a general method that extracts harmonic and anharmonic FCs
to arbitrary neighbour shell from a single set of displacements by
solving a linear least-squares problem with symmetry constraints.
Widely used implementations for the IFC3 include thirdorder.py
(distributed with ShengBTE) \cite{li_shengbte_2014} and phono3py
\cite{togo_first_2015, togo_implementation_2023}.  The
latter also provides the infrastructure for computing phonon
linewidths, thermal conductivity, and spectral functions from the
extracted IFC3.

\paragraph{Supercell convergence and aliasing.}
The choice of supercell determines the real-space range over which the
force constants can be resolved before periodic wrap-around introduces
aliasing.  For a supercell of linear dimension $N$ along a lattice
direction, the calculated force constant at separation $r$ is the
periodically folded sum
\begin{equation}\label{eq:fc_aliasing}
  \Phi_{\mathrm{per}}^{(N)}(r)
  = \sum_{p\in\mathbb{Z}} \Phi(r + pN).
\end{equation}
In a $2\times 2\times 2$ supercell, the resolvable relative
translations are limited to those representable within the finite
periodic cell -- contributions from more distant cells are folded back
onto these separations.  A larger supercell therefore allows force
constants to be resolved over a wider real-space range before aliasing
occurs \cite[Chs.~2,~6]{dove_introduction_1993}.

This aliasing directly affects the phonon dispersion. At wavevectors
$\vq$ that are not commensurate with the supercell mesh, the
Fourier-interpolated dynamical matrix inherits errors from the
folded-back tails of $\Phi$. In practice, convergence must therefore be
checked for the target observable---phonon frequencies, thermal
conductivity, or self-energy---by increasing the supercell size until
the aliased contributions are negligible
\cite{togo_first_2015, parlinski_first-principles_1997}.

For the IFC3, the aliasing issue is compounded by the fact that three
cell indices are involved: each pair in the triplet
$(0\kappa, l'\kappa', l''\kappa'')$ must lie within the
interaction range resolvable by the supercell.  This makes third-order
properties more demanding in supercell size than their harmonic
counterparts \cite{li_shengbte_2014, esfarjani_method_2008}.

% ----------------------------------------------------------------------
\subsubsection{Density-functional perturbation theory (DFPT)}
\label{ssub:dfpt_fc}
% ----------------------------------------------------------------------
Density functional perturbation theory avoids explicit supercell
calculations by computing the linear response of the electron density
to a phonon perturbation at an arbitrary wave-vector $\vq$ within the
primitive cell.  The method was introduced by
\citet{baroni_greens-function_1987} and
\citet{gonze_adiabatic_1995}, with a comprehensive review by
\citet{baroni_phonons_2001}.  Second-order force constants at any
$\vq$ are obtained directly from the self-consistent solution of the
Sternheimer equations, and real-space IFC2 are obtained by Fourier
interpolation of the dynamical matrix computed on a regular mesh of
$\vq$-points \cite{gonze_interatomic_1994}.

For cubic force constants, the $2n{+}1$ theorem allows $\Phi^{(3)}$
to be extracted from the first-order wavefunction and density
responses already computed at the IFC2 stage, without requiring
second-order responses \cite{debernardi_anharmonic_1995, debernardi_third-order_1994}.
\citet{deinzer_ab_2003} applied this
to compute phonon linewidths in Ge and Si, and
\citet{paulatto_anharmonic_2013} generalised it to arbitrary
triplets $(\vq_1,\vq_2,\vq_3)$ with $\vq_1{+}\vq_2{+}\vq_3 =\vc{G}$,
enabling the direct evaluation of three-phonon scattering matrix
elements throughout the Brillouin zone without real-space
supercells.  \citet{paulatto_first-principles_2015} later extended
the approach to compute phonon spectral functions with strong anharmonicity.

DFPT-based codes differ substantially in what they support at the
IFC2 and IFC3 levels.  For harmoinc IFCs, Quantum ESPRESSO
\cite{giannozzi_quantum_2009, giannozzi_advanced_2017} (via
\texttt{ph.x}) supports norm-conserving (NC), ultrasoft (USPP), and
projector-augmented-wave (PAW) pseudopotentials, with LDA, GGA, and
(in recent versions) hybrid functionals. ABINIT
\cite{gonze_abinit_2009} similarly supports both NC and PAW datasets,
non-magnetic, collinear, and non-collinear spin, as well as DFT$+U$.
VASP \cite{kresse_efficient_1996} provides a more rudimentary
DFPT implementation (\texttt{IBRION}\,{=}\,7,\,8) that is restricted
to displacements commensurate with the supercell ($\vq{=}\vc 0$) and
to LDA/GGA functionals.

For cubic force constants, the picture is more restrictive.  In
Quantum ESPRESSO, $\Phi^{(3)}$ is available through the D3Q and
\texttt{thermal2} packages
\cite{paulatto_anharmonic_2013, fugallo_ab_2013}.  These
codes support only norm-conserving pseudopotentials and only LDA or
GGA exchange--correlation functionals. Neither USPP, PAW, meta-GGA,
hybrid functionals, nor DFT$+U$ are currently implemented at third
order.  The same limitation applies to all third-order energy
derivatives in Quantum ESPRESSO, including Raman cross-sections.  In
ABINIT, a framework for third-order derivatives of the total energy
(3DTE) exists and is used in literature for non-linear optical
responses and Raman tensors \cite{veithen_nonlinear_2005, romero_abinit_2020}.
The extraction of the full three-phonon IFC3
tensor on a $\vq$-grid analogous to D3Q, however, is not available and, in
practice, anharmonic lattice dynamics in ABINIT is usually performed via
the supercell-based a-TDEP module \cite{bottin_-tdep_2020} rather than through pure
DFPT.  VASP does not currently implement third-order DFPT, and cubic
force constants with VASP are typically obtained through
finite--displacement calculations with, e.g., phono3py
\cite{togo_distributions_2015}. Hence, when computing IFC3 with DFPT, one is
effectively restricted to a LDA/GGA exchange correlation with
norm-conserving pseudopotentials, via Quantum ESPRESSO + D3Q.

% ----------------------------------------------------------------------
\subsubsection{Compressive sensing and machine-learning approaches}
\label{ssub:csld_fc}
% ----------------------------------------------------------------------
Compressive-sensing lattice dynamics (CSLD) casts the force-constant
extraction as an $\ell_1$-regularised regression on DFT forces
computed for a small number of randomly displaced configurations
\cite{zhou_lattice_2014, zhou_compressive_2019}.  The sparsity prior
selects the minimal set of non-negligible FC entries, potentially
capturing long-range couplings that would require prohibitively large
supercells in the standard finite-displacement scheme.  The hiPhive
package \cite{eriksson_hiphive_2019} implements this idea and uses
crystal-symmetry reduction and advanced regression
algorithms to extract force constants up to arbitrary order
from semi-random displacement configurations. These regression-based approaches are
particularly attractive when the full enumeration of finite
displacements becomes prohibitively expensive for large systems.

% ----------------------------------------------------------------------
\subsection{FC3 Compression}
\label{sub:fc3_compression}
% ----------------------------------------------------------------------

Storing the full FC3 tensor requires $\cO(\nprim\,\ndof^{2}) = \cO(\Ncell^{2}\nprim^{3})$
memory, which is prohibitive for large-scale simulations. Beyond storage, the dense four-tensor
contraction entering the phonon--phonon self-energy scales as $\cO(\ndof^{4})$ per frequency
point and dominates the cost of the quantum-transport simulation. It is hence desirable to
efficiently approximate the FC3 interactions.

Effectively every published anharmonic phonon NEGF implementation avoids the problem by
dropping self-energy and FC3 entries at large atomic separations
\cite{guo_quantum_2020, guo_anharmonic_2021, luisier_atomistic_2012, li_anharmonic_2009,
      polanco_nonequilibrium_2021, mingo_anharmonic_2006}.
The standard choice is a real-space pair-distance cut-off: keep only those triplets
$(0\kappa, l'\kappa', l''\kappa'')$ whose largest pairwise distance lies below some
$r_\mathrm{c}$. The same idea is used in the BTE workflows of ShengBTE and its successors,
which restrict FC3 to a few nearest-neighbour shells \cite{li_shengbte_2014}.

The drawback is that the approximation error is uncontrolled: there is no variational bound
on what is discarded, and the truncation removes genuine non-local contributions that can
matter quantitatively in the phonon self-energy. In anharmonic NEGF this is already visible in
silicon. \citeauthor{guo_quantum_2020} use first-principles FC3 in a quantum-transport treatment
of cross-plane heat flow in Si thin films, where they compare first- and third-nearest-neighbour
FC3 cutoffs \cite{guo_quantum_2020}. Using the same FC3 data in ShengBTE, they report bulk Si
thermal conductivities at 300~K of $120.69$, $136.46$, and $147.13~\mathrm{W\,m^{-1}K^{-1}}$ when
FC3 was truncated at the first, second, and third neighbour shell, respectively, so that a
first-shell cut-off underestimates the converged bulk value by nearly $20\%$. For bulk silicon,
including up to 4th nearest neighbour interactions gives converged results
\cite{li_shengbte_2014, qin_accelerating_2018}.

For three-dimensional covalent crystals, the cut-off problem is usually less severe than in
two-dimensional materials. \citeauthor{mcgaughey_phonon_2019} report heat conductivities
$\kappa_\ell(300\,\mathrm{K}) = 150, 152, 155, 151,$ and $149~\mathrm{W\,m^{-1}K^{-1}}$ for
cubic cut-offs of $2$, $3$, $4$, $5$, and $6$ neighbour shells, respectively in bulk Si, indicating
that once a few shells are retained the residual cut-off dependence is already small
\cite{mcgaughey_phonon_2019}. \citeauthor{jiang_how_2022} reached the same qualitative conclusion
by contrasting graphene and silicene against bulk Si: whereas the 2D systems require inclusion of
up to 14th order NN shells, bulk Si shows ``no convergence issue'' even for very small neighbour
cut-offs \cite{jiang_how_2022}.

That behaviour is, however, not universal across 3D materials. \citeauthor{qin_accelerating_2018}
showed that bulk SnSe still exhibits strong interactions at distances around $6.15~\text{\AA}$,
and that these coincide with a sharp change in the converged $\kappa_\ell$ once the pair cut-off
exceeds roughly $6~\text{\AA}$ \cite{qin_accelerating_2018}. They also show that for the 2D
materials studied, the observables in BTE diverge for increasing $q$-grid when the cut-off is too
small.

We therefore seek a representation of the FC3 tensor that
\begin{itemize}
    \item provides a controllable approximation error,
    \item achieves significant compression of the stored data, and
    \item turns the dense $\cO(\ndof^{4})$ self-energy contraction into a sum of low-order
          contractions whose cost scales polynomially in a parameter~$R$.
\end{itemize}
A large body of work has developed such representations in similar settings.
In quantum chemistry, tensor hypercontraction (THC) \cite{hohenstein_tensor_2012,
parrish_tensor_2012} and density fitting \cite{whitten_coulombic_1973} factorise the four-index
electron-repulsion integral tensor into products of matrices. In many-body physics, tensor
networks (matrix product states, tensor trains) \cite{oseledets_tensor-train_2011,
cichocki_tensor_2016} compress high-dimensional objects at controlled error. For the present
setting, lattice dynamics, tensor decomposition has recently proved effective:
\citeauthor{luo_data-driven_2024} used a truncated SVD on the electron-phonon matrix in Wannier
basis and recovered high-accuracy transport, spin, and superconducting observables from $1$--$2\%$
of the singular values \cite{luo_data-driven_2024}. The same group then proposed the permanent
CP (PCP) ansatz for the $n$-IFC tensors and reported $10^{3}$--$10^{4}$ compression factors
across Si, HgTe, MgO, TiNiSn and monoclinic ZrO$_2$ while retaining $>98\%$ accuracy on
$\kappa_\ell$ for three- and four-phonon scattering \cite{luo_tensor_2025}. Compressive-sensing
lattice dynamics (CSLD) is a method applied to a related problem at the fitting stage: an
$\ell_1$-regularised regression on DFT forces selects a sparse set of non-negligible FC entries
directly \cite{zhou_lattice_2014, zhou_compressive_2019}. Projector-based approaches, implemented
in the symfc software, construct a complete orthonormal basis for FCs compatible with all
crystal symmetries \cite{seko_projector-based_2024}, while group-theoretical reductions enumerate
the independent FC3 entries from the space group \cite{fu_group_2019}. In the following, we take
the fully-populated FC3 tensor as given (from finite displacements, DFPT, or CSLD) and compress
its representation. However, one can likely also fit a compressed tensor from DFT data with an
adaption of CSLD or similar methods without ever obtaining a full FC3 tensor.

\citet{kolda_tensor_2009} organise tensor compression into two principal families: the
Tucker decomposition, which writes a tensor as a core tensor multiplied by a factor matrix
along each mode (generalising principal-component analysis), and the CP
(CANDECOMP/PARAFAC) decomposition, which writes it as a sum of rank-one outer products
(generalising the matrix SVD in a different direction). Applied to a tensor with the FC3 index
pattern and combined with the possible symmetry constraints on the factor matrices, these two
families yield five canonical ans\"atze, ordered here from least to most
tensor-structure-aware:
\begin{enumerate}
  \item Truncated matricization SVD, i.e.\ the classical matrix SVD applied to an matricization of
        the tensor. fully algebraic and Eckart--Young optimal in the chosen matricization, but
        separating only one leg from the other two.
  \item Truncated HOSVD (Tucker), with independent factor matrices per mode and a small dense core.
  \item CP, obtained by restricting the Tucker core to be diagonal, with three
        independent rank-$R$ factor matrices and no symmetry constraint.
  \item INDSCAL \cite{carroll_analysis_1970}, obtained from CP by setting the
        factor matrices on the two internal legs equal, exploiting the
        $S_2$ symmetry $\Phi_{\mu j k} = \Phi_{\mu k j}$.
  \item Symmetric CP, equivalent to the Waring decomposition of a homogeneous
        cubic polynomial \cite{brachat_symmetric_2010, comon_symmetric_2008,
        oeding_eigenvectors_2013}, obtained from CP by setting all three factor
        matrices equal (using the $p(\mu)$ lift of a primitive-cell DOF to its
        supercell counterpart in cell zero), exploiting the full $S_3$
        symmetry.
\end{enumerate}
Entries 1 and 2 are non-iterative. Entries 3--5 are NP-hard in the worst case
\cite{hillar_most_2013} and use iterative optimisation, with the known caveats of
rank-degeneracy and non-closedness of rank-$R$ sets
\cite{stegeman_degeneracy_2006, de_silva_tensor_2008}. Each added symmetry constraint from
entry~2 onwards reduces the parameter count at matched rank while shrinking the set of
representable tensors. As discussed in %\cref{sub:fc3_observable_error}
\todo{write section about relation between error and observable}
, the most-constrained ansatz is not necessarily the least accurate per
parameter.

% ----------------------------------------------------------------------
\subsubsection{Truncated matricization SVD}
\label{sub:fc3_matsvd}
% ----------------------------------------------------------------------
Stack the $3\nprim$ slices $M_\mu = \Phi_{\mu,:,:}\in\mathbb{R}^{\ndof\times\ndof}$
into a single matrix
\begin{equation}\label{eq:fc3_matsvd_stack}
  M_{\mathrm{stack}}
  = \begin{pmatrix} M_1 \\ \vdots \\ M_{3\nprim} \end{pmatrix}
  \in \mathbb{R}^{3\nprim\,\ndof \,\times\, \ndof}.
\end{equation}
A truncated matrix SVD at rank $R$ gives
\begin{equation}\label{eq:fc3_matsvd_ansatz}
    \tilde\Phi_{\mu j k}^{\mathrm{mSVD}}
    = \sum_{r=1}^{R} F_r(\mu,j)\; h_r(k),
\end{equation}
with $F_r(\mu,j) = \sigma_r\,[\vc{u}_r]_{(\mu-1)\ndof+j}$ and $h_r(k)=[\vc{v}_r]_k$, where
$\sigma_r$, $\vc{u}_r$, $\vc{v}_r$ are the $R$ leading singular values, left singular vectors,
and right singular vectors of $M_{\mathrm{stack}}$. This is equivalent to a Tucker1 decomposition
with compression on mode~3 only \cite[\S4]{kolda_tensor_2009}.

\paragraph{Computation.} A single matrix SVD of $M_{\mathrm{stack}}$, computed by
standard dense LAPACK routines or, for large $\ndof$, by randomised linear algebra algorithms
\cite{halko_finding_2011}. The Eckart--Young--Mirsky theorem
\cite{eckart_approximation_1936, mirsky_symmetric_1960} guarantees that, at the chosen $(\mu,j)|k$ bipartition,
the truncation is globally optimal among all rank-$R$ approximations in the Frobenius norm.
This makes it the natural baseline against which the structurally richer decompositions of
\cref{sub:fc3_tucker,sub:fc3_cp,sub:fc3_indscal,sub:fc3_symcp} must be compared. The ansatz does
not exploit the $S_2$ symmetry on $(j,k)$: modes~1 and~2 are folded into a single flat index, so
parameters are wasted on storing redundant information and the reconstruction is not exactly
symmetric in $(j,k)$ unless restored in post-processing.

\paragraph{Parameter count.}
\begin{equation}\label{eq:fc3_matsvd_params}
  N_{\mathrm{param}}^{\mathrm{mSVD}}(R)
  = R\bigl((3\nprim+1)\,\ndof + 1\bigr)
  \;\sim\; \cO\!\bigl(R\,\nprim\,\ndof\bigr)
  \;=\; \cO\!\bigl(R\,\Ncell\,\nprim^{2}\bigr),
\end{equation}
one factor of $\nprim$ worse than the CP-family ans\"atze below.

\paragraph{Self-energy kernel.} The separable structure factorises the ring contraction entering
the phonon--phonon self-energy into two independent bilinear forms. The $\vq$-convolution is an
FFT over the $R^{2}$ rank pairs, replacing the dense $\ndof^{4}$ contraction by
$R^{2}\ndof$-scaling accumulations.

% ----------------------------------------------------------------------
\subsubsection{Truncated HOSVD (Tucker decomposition)}
\label{sub:fc3_tucker}
% ----------------------------------------------------------------------
The Tucker ansatz writes FC3 as a multilinear contraction of a small core tensor with one factor
matrix per mode \cite{tucker_mathematical_1966, kolda_tensor_2009}:
\begin{equation}\label{eq:fc3_tucker_ansatz}
    \tilde\Phi_{\mu j k}^{\mathrm{HOSVD}}
    = \sum_{p=1}^{R_1}\sum_{q=1}^{R_2}\sum_{r=1}^{R_2}
      G_{p q r}\; A_{\mu p}\; B_{j q}\; B_{k r},
\end{equation}
with $A\in\mathbb{R}^{3\nprim\times R_1}$, $B\in\mathbb{R}^{\ndof\times R_2}$ columnwise
orthonormal, and a core $G\in\mathbb{R}^{R_1\times R_2\times R_2}$. To preserve the $S_2$
symmetry of the two internal legs, the factor matrices on modes~2 and~3 can be chosen identical ($B=C$)
and the core is constrained to satisfy $G_{pqr}=G_{prq}$.

\paragraph{Computation.} The truncated higher-order SVD (HOSVD) of
\citet{de_lathauwer_multilinear_2000} produces $A$ and $B$ in closed form as the $R_1$ and $R_2$
leading left singular vectors of the mode-$1$ and mode-$2$ unfoldings of $\Phi$, followed by the
core contraction $G = \Phi\times_1 A^\top\times_2 B^\top\times_3 B^\top$ (which is automatically
$S_2$-symmetric when $\Phi$ is). This is quasi-optimal: the truncation error is bounded by a
factor $\sqrt{3}$ above the best rank-$(R_1,R_2,R_2)$ approximation error
\cite{hackbusch_tensor_2019, oseledets_tensor-train_2011}. The fit can optionally be refined by a
few steps of higher-order orthogonal iteration (HOOI), an alternating least squares (ALS)-type loop
 \cite{de_lathauwer_best_2000}. The special case $R_1=3\nprim$, $A=I$ (compression of the two
internal legs only) is the Tucker1 baseline and is obtained by a single SVD of the corresponding
unfolding.

\paragraph{Parameter count.}
\begin{equation}\label{eq:fc3_tucker_params}
  N_{\mathrm{param}}^{\mathrm{HOSVD}}(R_1,R_2)
  = 3\nprim\,R_1 + \ndof\,R_2 + R_1\,\tfrac{R_2(R_2+1)}{2}
  \;\sim\; \cO\!\bigl(R_2\,\ndof + R_1\,R_2^{2}\bigr),
\end{equation}
where the $S_2$ constraint halves the nominal core cost $R_1 R_2^{2}$ to $R_1 R_2(R_2+1)/2$.
The $R_2^{2}$-scaling core term makes Tucker the most parameter-heavy of the five ans\"atze
when $R_2$ is not small compared to $\ndof^{1/2}$. for strongly anharmonic systems where the
effective rank is moderate this is mitigated by allowing $R_1\ll R_2$.

\paragraph{Self-energy kernel.} The self-energy contraction factorises into three mode-wise
contractions against $A$ and $B$, leaving the $R_1\times R_2\times R_2$ core as a small dense
sub-problem. The total cost scales as $R_2^{2}(R_1+\ndof)$ per frequency point rather than
$\ndof^{4}$.

% ----------------------------------------------------------------------
\subsubsection{Canonical Polyadic decomposition (CP)}
\label{sub:fc3_cp}
% ----------------------------------------------------------------------
Dropping the off-diagonal entries of the Tucker core yields the CP approximation, also
known as CANDECOMP and PARAFAC \cite{harshman_foundations_1970, hitchcock_expression_1927, carroll_analysis_1970}, a sum of
rank-one outer products with three independent factor vectors per term:
\begin{equation}\label{eq:fc3_cp_ansatz}
    \tilde\Phi_{\mu j k}^{\mathrm{CP}}
    = \sum_{r=1}^{R} \lambda_r\;a_r(\mu)\;b_r(j)\;c_r(k).
\end{equation}
No symmetry between the modes is imposed. For a tensor that is
$S_2$- or $S_3$-symmetric, the fit does not necessarily converge to a symmetric
solution \cite{ten_berge_typical_2004}. Uniqueness can be shown under certain conditions
\cite{kruskal_three-way_1977, sidiropoulos_uniqueness_2000} and is generically
guaranteed for the FC3 tensor of interest here
\cite{chiantini_generic_2012}.

\paragraph{Computation.} The standard algorithm is alternating least squares
(CP-ALS), which cycles through the three modes solving, at each step, a
linear least-squares problem for one factor matrix while the other two are
fixed \cite[Sec.~3.4]{kolda_tensor_2009}. Convergence is only to a
stationary point and can be slow. Variants include
damped Gauss--Newton (dGN), PMF3 \cite{paatero_weighted_1997},
line-search extrapolation (ELS) \cite{rajih_enhanced_2008}, and
simultaneous-diagonalisation approaches \cite{de_lathauwer_link_2006}
that recast CP as a generalised eigenproblem when
$R\le\min(3\nprim,\ndof)$.

\paragraph{Parameter count.}
\begin{equation}\label{eq:fc3_cp_params}
  N_{\mathrm{param}}^{\mathrm{CP}}(R)
  = R\,\bigl(3\nprim + 2\,\ndof + 1\bigr)
  \;\sim\; \cO\!\bigl(R\,\ndof\bigr)
  \;=\; \cO\!\bigl(R\,\Ncell\,\nprim\bigr).
\end{equation}

% ----------------------------------------------------------------------
\subsubsection{INDSCAL (CP with internal-leg symmetry)}
\label{sub:fc3_indscal}
% ----------------------------------------------------------------------
Individual-differences-in-scaling (INDSCAL) \cite{carroll_analysis_1970, kolda_tensor_2009} is the
specialisation of CP in which the factor matrices on two modes are constrained equal, tailored
to tensors whose slices $\Phi_{\mu,:,:}$ are symmetric:
\begin{equation}\label{eq:fc3_indscal_ansatz}
    \tilde\Phi_{\mu j k}^{\mathrm{INDSCAL}}
    = \sum_{r=1}^{R} d_r(\mu)\; v_r(j)\; v_r(k),
\end{equation}
with $d_r\in\mathbb{R}^{3\nprim}$ and $v_r\in\mathbb{R}^{\ndof}$. This imposes the $S_2$
symmetry $\Phi_{\mu j k}=\Phi_{\mu k j}$. Typical and maximal ranks for
partially symmetric tensors were characterised by
\citet{ten_berge_typical_2004}: for the shapes relevant here, they coincide with or
are strictly smaller than the ranks of generic nonsymmetric tensors of the same shape, so the
symmetry constraint is not costly in expressive power.

\paragraph{Computation.} \citet{carroll_analysis_1970} originally proposed running unconstrained
CP-ALS with the two internal factors $B$ and $C$ updated separately, relying on symmetry of the
data to drive $B\to C$. \citet{ten_berge_typical_2004, ten_berge_carroll_2009} showed that this convergence is generic
but not guaranteed. Modern implementations impose equality explicitly via a symmetric ALS or via
a Gauss--Newton method \cite{kolda_tensor_2009, kolda_numerical_2015}.
A cheap closed-form alternative performs a matrix SVD of $\Phi$ reshaped to
$3\nprim\times\ndof^{2}$ followed by an eigendecomposition of each symmetric slice of the
resulting right factors, and keeps the $R$ terms with largest amplitude. This is not jointly
optimal (the total rank needed to match a given error is typically larger than that of a
jointly-fitted INDSCAL) but is useful as an initialisation.

\paragraph{Parameter count.}
\begin{equation}\label{eq:fc3_indscal_params}
  N_{\mathrm{param}}^{\mathrm{INDSCAL}}(R)
  = R\,\bigl(3\nprim + \ndof\bigr)
  \;\sim\; \cO\!\bigl(R\,\ndof\bigr)
  \;=\; \cO\!\bigl(R\,\Ncell\,\nprim\bigr),
\end{equation}
one factor of $\ndof$ smaller than CP at matched rank, from merging $b_r\equiv c_r$.

\paragraph{Self-energy kernel.} Identical in scaling to CP, $\cO(R^{2}\ndof)$ per frequency, but
with the additional property that the reconstruction is exactly $S_2$-symmetric in $(j,k)$, so
ASR-sensitive observables receive only the quadratic-in-error contribution discussed in
\todo{actually discuss this}%\cref{sub:fc3_observable_error}.


% ----------------------------------------------------------------------
\subsubsection{Symmetric CP (Waring) decomposition}
\label{sub:fc3_symcp}
% ----------------------------------------------------------------------
The symmetric CP decomposition imposes equality of all three factor matrices, $A=B=C$, on the
fully $S_3$-symmetrised FC3 tensor obtained via the supercell lift
$p:\{1,\dots,3\nprim\}\to\{1,\dots,\ndof\}$. The Waring parametrisation reads:
\begin{equation}\label{eq:fc3_waring_ansatz}
    \tilde\Phi_{\mu j k}^{\mathrm{Waring}}
    = \sum_{r=1}^{R} \lambda_r\;
      v_r\!\bigl(p(\mu)\bigr)\; v_r(j)\; v_r(k),
\end{equation}
with $v_r\in\mathbb{R}^{\ndof}$ and $\lambda_r\in\mathbb{R}$. This is the tensor form of the
classical Waring problem: expressing a homogeneous cubic polynomial as a sum of $R$ cubes of
linear forms \cite{brachat_symmetric_2010, comon_symmetric_2008, oeding_eigenvectors_2013}. For
generic third-order symmetric tensors on $\ndof$ indices, the Waring rank equals
$\lceil\binom{\ndof+2}{3}/\ndof\rceil$ by the Alexander--Hirschowitz theorem
\cite{alexander_generic_1997} (which is a non-trivial statement and proven by Alexander and Hirschowitz
in a set of 5 papers of over 100 pages). An alternative ``permanent-CP'' parametrisation, introduced
by \citet{luo_tensor_2025} as PCP, uses three distinct vectors per term explicitly symmetrised
over $S_3$,
\begin{equation}\label{eq:fc3_pcp_ansatz}
    \tilde\Phi_{\mu j k}^{\mathrm{PCP}}
    = \sum_{\xi=1}^{R} \frac{\lambda_\xi}{6}
      \sum_{\sigma\in S_3}
      u^{\xi}_{\sigma_1}\!\bigl(p(\mu)\bigr)\;
      u^{\xi}_{\sigma_2}(j)\;
      u^{\xi}_{\sigma_3}(k),
\end{equation}
with $u^{\xi}_1, u^{\xi}_2, u^{\xi}_3\in\mathbb{R}^{\ndof}$. PCP fits the same class of symmetric
tensors as Waring but with more capacity per term (at threefold parameter cost) and a smoother
optimisation landscape.

\paragraph{Computation.} Algorithms fall into two categories. Algebraic methods reconstruct the
decomposition from moment and flattening matrices via generalised eigenproblems and are exact up
to numerical precision when the rank is below the generic bound
\cite{brachat_symmetric_2010, oeding_eigenvectors_2013}. Numerical methods optimise
$\|\Phi-\tilde\Phi\|_F^{2}$ directly: the shifted symmetric higher-order power method (SS-HOPM)
\cite{kolda_shifted_2011} and the subspace power method \cite{kileel_subspace_2025} extract
components one at a time with provable convergence for low-rank tensors; all-at-once gradient
and Gauss--Newton methods \cite{kolda_numerical_2015} fit all $R$ components simultaneously.
Degeneracy and diverging-norm instabilities can appear because the set of symmetric tensors of
rank $\le R$ is not closed in general
\cite{stegeman_degeneracy_2006, de_silva_tensor_2008}. Remedies include norm penalties and
second-order line searches.

\paragraph{Parameter count.}
\begin{align}
  N_{\mathrm{param}}^{\mathrm{Waring}}(R) &= R\,(\ndof + 1)
    \;\sim\; \cO\!\bigl(R\,\ndof\bigr)
    \;=\; \cO\!\bigl(R\,\Ncell\,\nprim\bigr),
    \label{eq:fc3_waring_params}\\
  N_{\mathrm{param}}^{\mathrm{PCP}}(R) &= R\,(3\,\ndof + 1)
    \;\sim\; \cO\!\bigl(R\,\ndof\bigr)
    \;=\; \cO\!\bigl(R\,\Ncell\,\nprim\bigr).
    \label{eq:fc3_pcp_params}
\end{align}

\paragraph{Self-energy kernel.} Identical in scaling to CP and INDSCAL, $\cO(R^{2}\ndof)$ per
frequency, but with the additional property that the reconstruction is exactly $S_3$-symmetric
in $(p(\mu),j,k)$, so ASR-sensitive observables receive only the quadratic-in-error contribution
discussed in \todo{discuss}
%\cref{sub:fc3_observable_error}.


% ----------------------------------------------------------------------
\subsubsection{Further approximations}
\label{sub:fc3_other}
% ----------------------------------------------------------------------
The five ans\"atze of
\cref{sub:fc3_matsvd,sub:fc3_tucker,sub:fc3_cp,sub:fc3_indscal,sub:fc3_symcp} do not exhaust the
useful choices. Further approximations appear in the broader literature and can either be
combined with them or replace them in specific regimes.

\paragraph{Cross / CUR / skeleton approximations.} Rather than fitting a compressed form to a
precomputed $\Phi$, one can build the approximation from a small number of sampled fibers of the
tensor and interpolate the rest:
\begin{equation}
  \tilde\Phi_{\mu j k}^{\mathrm{CUR}}
  = \sum_{p,q,r}\Phi_{\mu,J_q,K_r}\,
    \Phi_{I_p,J_q,K_r}^{-1}\,\Phi_{I_p,j,k},
\end{equation}
where $I,J,K$ are index subsets of size $\sim R$ selected to maximise the volume of the
intersection subtensor, and $\Phi_{I,J,K}$ is inverted in the appropriate pseudoinverse sense.
This bypasses the construction of the full FC3 altogether, which is normally the computational
bottleneck (finite-displacement DFT). Rank-$R$ tensor cross
\cite{oseledets_tt-cross_2010, savostyanov_quasioptimality_2014} accesses $\cO(R^{2}\ndof)$ entries.
Adaptive cross approximation \cite{bebendorf_approximation_2000} and CUR-style subset selection
\cite{goreinov_theory_1997, mahoney_cur_2009} give analogous matrix-level
guarantees, and fiber-sampling Tucker (FSTD) of \citet{caiafa_generalizing_2010} yields an explicit
Tucker form from $\cO(R^{N-1})$ fibers. The tensor-CURT extension of
\citet{song_relative_2019} provides the first relative-error bounds in this setting. For
FC3 the main challenge is that cross methods have no intrinsic notion of ASR or $S_3$ symmetry.
A hybrid that samples in the symmetric CP parameter space would combine the two sets of
guarantees.

\paragraph{Tensor train and hierarchical Tucker.} The tensor-train (TT) format
\cite{oseledets_tensor-train_2011} writes an $N$-index tensor as a contraction of $N$ three-index
cores with bond ranks $(R_1,\dots,R_{N-1})$. For third-order tensors this collapses to
\begin{equation}
  \tilde\Phi_{\mu j k}^{\mathrm{TT}}
  = \sum_{\alpha_1,\alpha_2} G_1(\mu,\alpha_1)\,
    G_2(\alpha_1,j,\alpha_2)\,G_3(\alpha_2,k),
\end{equation}
which is structurally a Tucker1 SVD of the $(\mu)|(j,k)$ unfolding followed by a Tucker1 SVD of
the $(\alpha_1 j)|(k)$ matricisation of the residual. It does not offer independent advantages
beyond \cref{sub:fc3_tucker}. TT and the hierarchical Tucker (HT) format of
\citet{hackbusch_tensor_2019} become first-class methods at order
$\ge 4$ and are the natural representation for FC4 (four-phonon) tensors, where CP and Tucker
lose efficiency. \citet{lee_solving_2026} recently used Matrix Product States (= TT) to solve the
Peierls--Boltzmann equation in silicon.

\paragraph{Block-term decomposition (BTD).} \citet{de_lathauwer_decompositions_2008, de_lathauwer_decompositions_2008-1}
generalise both CP and Tucker by writing the tensor as a sum of low-multilinear-rank Tucker blocks,
\begin{equation}
  \tilde\Phi_{\mu j k}^{\mathrm{BTD}}
  = \sum_{r=1}^{R}\bigl(\mathcal{G}_r\times_1 A_r\times_2 B_r\times_3 C_r\bigr)_{\mu j k},
\end{equation}
with each $\mathcal{G}_r$ of small core shape (e.g.\ the rank-$(L,L,1)$ BTD has
$\mathcal{G}_r$ of shape $L\times L\times 1$). CP is the rank-$(1,1,1)$ special case; Tucker is
the single-block special case. For FC3 this structurally expresses block patterns like
per-atom-species or per-shell contributions in a single decomposition, and inherits the
uniqueness and ALS algorithms of the parent families.

\paragraph{Constrained / symmetry-adapted decompositions (CANDELINC).} CANDELINC
\cite{carroll_candelinc_1980} is CP with factor matrices constrained to user-chosen
linear subspaces,
\begin{equation}
  A=\Phi_A\hat A,\quad B=\Phi_B\hat B,\quad C=\Phi_C\hat C,
\end{equation}
where $\Phi_A,\Phi_B,\Phi_C$ are columnwise-orthonormal basis matrices. It reduces to an ordinary
CP fit in the reduced coordinates $(\hat A,\hat B,\hat C)$ applied to the projected tensor
$\hat\Phi = \Phi\times_1\Phi_A^\top\times_2\Phi_B^\top\times_3\Phi_C^\top$. For FC3 this is the
natural framework to couple any of the five ans\"atze with a space-group-compatible basis:
choose $\Phi_A,\Phi_B,\Phi_C$ to span the ASR-respecting, symmetry-invariant subspace (symfc
\cite{seko_projector-based_2024}, Fu--Marianetti \cite{fu_group_2019}, CSLD
\cite{zhou_lattice_2014, zhou_compressive_2019}), then fit INDSCAL or symmetric CP on the reduced
tensor.

\paragraph{Sparse / structured CP.} Imposing sparsity or smoothness priors on the CP factor
vectors yields decompositions whose rank-$R$ components are localised on a few atoms or DOFs.
As an example, one can penalise the loss with
\begin{equation}
  L(A,B,C) = \tfrac12\|\Phi - ( A,B,C)\|_{F}^{2}
            + \lambda_1\sum_r\bigl(\|a_r\|_1+\|b_r\|_1+\|c_r\|_1\bigr),
\end{equation}
optionally replacing $\|\cdot\|_1$ by group-sparse or total-variation terms
\cite{papalexakis_parcube_2012, kim_sparse_2007}, and are solved by
ALS or AO-ADMM.

% ----------------------------------------------------------------------
\section{Phonon NEGF}
\label{sec:ph_negf}
% ----------------------------------------------------------------------

We apply the non-equilibrium Green's function (NEGF) formalism to phonon quantum transport,
taking as starting point the energy expansion in atomic displacements introduced in
\cref{eq:e_exp_ph}. Truncated at second order, the Hamiltonian generates
ballistic (harmonic) transport. The third-order contribution gives the
leading anharmonic phonon--phonon scattering, treated perturbatively in
the self-consistent Born approximation (SCBA). The phonon NEGF algorithm is
in close analogy with the electronic case~\cite{caroli_direct_1971,meir_landauer_1992,haug_quantum_2008,wang_quantum_2008, wang_nonequilibrium_2014},
with two major differences: the natural energy variable is
$\omega^2$ because the lattice equation of motion is second-order in
time, and the particles are bosons whose vacuum-fluctuation
contributions persist at $T = 0$. The phonon-NEGF strategy was developed
by Mingo and Yang for ballistic transport~\cite{mingo_calculation_2003, mingo_phonon_2003},
extended to anharmonicity by~\citet{mingo_anharmonic_2006, wang_nonequilibrium_2007,wang_nonequilibrium_2007},
and by~\cite{luisier_atomistic_2012,lee_anharmonic_2018}. An explicit derivation
of the SCBA self-energy at the diagrammatic level is given by~\citet{guo_quantum_2020}.

% ----------------------------------------------------------------------
\subsection{Hamiltonian and lattice-dynamical setup}
\label{sub:ph_hamiltonian}
% ----------------------------------------------------------------------

Let $u_{l\kappa\alpha}$ denote the Cartesian displacement of atom $\kappa$
in unit cell $l$, and let $m_\kappa$ be its mass. In
mass-weighted coordinates $\tilde{u}_{l\kappa\alpha}=\sqrt{m_\kappa}\,u_{l\kappa\alpha}$,
the lattice Hamiltonian truncated at third order is
\begin{equation}\label{eq:H_phonon}
  \hat{H}
  = \tfrac{1}{2}\sum_{\mu}\tilde{p}_\mu^{\,2}
    + \tfrac{1}{2}\sum_{\mu\nu} D_{\mu\nu}\,\tilde{u}_\mu \tilde{u}_\nu
    + \tfrac{1}{3!}\sum_{\mu\nu\lambda} \Phi_{\mu\nu\lambda}\,
      \tilde{u}_\mu \tilde{u}_\nu \tilde{u}_\lambda
    + \mathcal{O}(\tilde{u}^4),
\end{equation}
where the multi-index $\mu = (l,\kappa,\alpha)$ collects unit-cell,
sublattice, and Cartesian labels. The mass-weighted second- and
third-order force constants are
\begin{align}
  D_{\mu\nu} &= \frac{1}{\sqrt{m_\kappa m_{\kappa'}}}\,\Phi^{(2)}_{\mu\nu}, &
  \Phi_{\mu\nu\lambda} &= \frac{1}{\sqrt{m_\kappa m_{\kappa'} m_{\kappa''}}}\,\Phi^{(3)}_{\mu\nu\lambda},
\end{align}
with $\Phi^{(2)}_{\mu\nu}=\partial^2 V/\partial u_\mu \partial u_\nu$ and analogously
for $\Phi^{(3)}$. The acoustic sum rules
\begin{equation}\label{eq:asr}
  \sum_{l\kappa} \Phi^{(n)}_{l\kappa\alpha,\,\mu_2\ldots\mu_n}
  = 0
  \qquad (n = 2, 3, \ldots),
\end{equation}
plus translational and permutation symmetry are imposed by the definition of the FCs.

Quantizing in the harmonic basis,
\begin{equation}
  \tilde{u}_\mu(t)
  = \sum_{s} \sqrt{\frac{\hbar}{2\omega_s}}\,e_{\mu s}\,
    \bigl[a_s\,e^{-i\omega_s t} + a_s^\dagger\,e^{i\omega_s t}\bigr],
  \qquad D\,e_s = \omega_s^2\,e_s,
\end{equation}
we obtain $\omega^2$ as the spectral variable of phonon Green's
functions.

% ----------------------------------------------------------------------
\subsection{Contour-ordered Green's function and Dyson equation}
\label{sub:ph_dyson}
% ----------------------------------------------------------------------

Following the conventions of~\citet{wang_nonequilibrium_2014}, we
define the contour-ordered phonon Green's function on the Keldysh contour $\mathcal{C}$:
\begin{equation}\label{eq:G_def}
  G_{\mu\nu}(\tau,\tau')
  = -\frac{i}{\hbar}\,
    \bigl\langle T_{\mathcal{C}}\,
      \tilde{u}_\mu(\tau)\,\tilde{u}_\nu(\tau')\bigr\rangle.
\end{equation}
Projecting $\tau$ onto the upper/lower branches yields the
lesser, greater, retarded, advanced, and time-ordered components
$G^{\lessgtr},\, G^{R/A},\, G^t$, with the standard relations
\begin{align}
  G^R(t,t') &= \theta(t-t')\bigl[G^>(t,t') - G^<(t,t')\bigr], &
  G^A &= (G^R)^\dagger,\\
  G^R - G^A &= G^> - G^<. &
  G^t &= G^< + G^r
\end{align}
At equilibrium, the fluctuation--dissipation theorem can be written as~\cite[Eq. 24]{wang_nonequilibrium_2014}
\begin{equation}\label{eq:phonon_fdt}
  G^<(\omega) = -2 i\,n_{\mathrm{BE}}(\omega,T)\,\mathrm{Im}\,G^R(\omega),\qquad
  G^>(\omega) = -2 i\,[1 + n_{\mathrm{BE}}(\omega,T)]\,\mathrm{Im}\,G^R(\omega).
\end{equation}
Out of
equilibrium, \cref{eq:phonon_fdt} does not hold and one must solve the
Keldysh equation
\begin{equation}\label{eq:keldysh_eq}
  G^{\lessgtr}(\omega) = G^R(\omega)\,\Sigma^{\lessgtr}(\omega)\,G^A(\omega).
\end{equation}
The retarded Green's function can be computed from the Dyson equation
\begin{equation}\label{eq:system}
  \bigl[G^R(\omega^2)\bigr]^{-1}
  = (\omega^2 + i\eta)\,\mathbb{I}
    - D
    - \Sigma^R(\omega^2),
\end{equation}
with $D$ the dynamical matrix, $\eta>0$ a numerical broadening matched
to the frequency grid, and the total retarded self-energy decomposed as
\begin{equation}
  \Sigma^R(\omega^2) = \Sigma_L^R(\omega^2) + \Sigma_R^R(\omega^2)
                      + \Sigma_s^R(\omega^2).
\end{equation}
The lead self-energies $\Sigma_{L,R}^R$ are due to the semi-infinite harmonic
contacts. $\Sigma_s^R$ collects all many-body scattering and equals zero
in the ballistic limit. The spectral function is defined as
\begin{equation}\label{eq:spectral}
  A(\omega^2) = i[G^R(\omega^2)-G^A(\omega^2)] = -2\,\mathrm{Im}\,G^R(\omega^2),
\end{equation}
with sum rule $\int_0^\infty\!\frac{d(\omega^2)}{2\pi}\,[A(\omega^2)]_{\mu\mu}
= 1$ and poles only at positive $\omega^2$.

% ----------------------------------------------------------------------
\subsection{Transverse periodic NEGF}
\label{sub:transper_negf}
% ----------------------------------------------------------------------

We consider both finite and transversely periodic geometries, with
transport along $x$ and periodicity in $(y,z)$. Following~\cite{tian_enhancing_2012,zhang_simulation_2007,guo_quantum_2020},
the transverse wave vector
$\qp$ (we use $q$ for phonons, $k$ for electrons) is sampled on a uniform
Monkhorst--Pack mesh
\begin{equation}
  k_i = \frac{n_i + 0.5}{N_i} - 0.5 + \mathrm{shift}_i,
  \qquad n_i = 0,\ldots,N_i-1,
\end{equation}
with total count $N_q$ chosen to converge the transverse Brillouin-zone
integral (e.g.\ $N_q=400$ on a $20\times 20$ grid). The Dyson equation
becomes $\qp$-dependent,
\begin{equation}\label{eq:dyson_qp}
  \bigl[G^R(\omega^2,\qp)\bigr]^{-1}
  = (\omega^2 + i\eta)\,\mathbb{I}
    - D(\qp)
    - \Sigma^R(\omega^2,\qp).
\end{equation}
For fixed $\qp = (q_x,q_y)$, the on-site and nearest-neighbour
transport-direction blocks are obtained by partial Fourier transform
over the $q_z$ mesh:
\begin{align}
  H_{00}(\qp)
    &= \frac{C}{N_{q_z}}
       \sum_{q=0}^{N_{q_z}-1}
       D^{(\mathrm{B})}\!\left(q_x,q_y,\frac{q}{N_{q_z}}\right),
  \label{eq:H00}\\
  H_{01}(\qp)
    &= \frac{C}{N_{q_z}}
       \sum_{q=0}^{N_{q_z}-1}
       D^{(\mathrm{B})}\!\left(q_x,q_y,\frac{q}{N_{q_z}}\right)\,
       e^{-2\pi i q / N_{q_z}}.
  \label{eq:H01}
\end{align}
At nonzero $\qp$, these blocks are generally complex and may equivalently
be written as the 2D Fourier transform of the real-space FC2:
\begin{equation}
  H_{0n_z}(\qp)
  = \sum_{n_x,n_y}
    \Phi(n_x,n_y,n_z)\,e^{2\pi i(q_x n_x + q_y n_y)},
\end{equation}
with hermiticity enforcing $H_{10}(\qp) = H_{01}^\dagger(\qp)$.

% ----------------------------------------------------------------------
\subsection{Contact self-energy}
\label{sub:contact_se}
% ----------------------------------------------------------------------

The contact self-energies $\Sigma^R_{L,R}$ embed the semi-infinite
harmonic leads via the surface Green's function $g_s$, computed
iteratively by the renormalization--decimation algorithm of~\citet{sancho_highly_1985}.
The method has quadratic-doubling convergence: each iteration
doubles the number of decimated layers, so the iteration error
falls as $\|T\|^{2^n}$ with $n$ the iteration index.

\paragraph{Decimation iteration.}
Given the bulk on-site and coupling blocks $H_{00}$ and $H_{01}$,
initialize
\begin{equation}
  \epsilon = \epsilon_s = (\omega^2 + i\eta)\mathbb{I} - H_{00},
  \qquad
  \alpha = -H_{10},
  \qquad
  \beta  = -H_{01},
\end{equation}
and iterate
\begin{align}
  \epsilon_s &\leftarrow \epsilon_s - \alpha\,\epsilon^{-1}\beta, \label{eq:sr_eps_s}\\
  \epsilon   &\leftarrow \epsilon - \alpha\,\epsilon^{-1}\beta
              - \beta\,\epsilon^{-1}\alpha, \label{eq:sr_eps}\\
  \alpha     &\leftarrow \alpha\,\epsilon^{-1}\alpha, \label{eq:sr_alpha}\\
  \beta      &\leftarrow \beta\,\epsilon^{-1}\beta, \label{eq:sr_beta}
\end{align}
terminating when $\|\alpha\|+\|\beta\|<\mathrm{tol}$. The surface
Green's function is then $g_s = \epsilon_s^{-1}$. For the left lead
$g_L = g_s(H_{00},H_{01})$. For the right lead, the coupling direction
is reversed, $g_R = g_s(H_{00},H_{01}^\dagger)$. The contact
self-energies follow from
\begin{equation}\label{eq:sigma_lr}
  \Sigma_L^R = D_{10}\,g_L\,D_{01},
  \qquad
  \Sigma_R^R = D_{01}\,g_R\,D_{10},
\end{equation}
and the broadening matrices, which set the rate at which lead $\ell$
injects phonons into the device, are
\begin{equation}\label{eq:gamma_def}
  \Gamma_\ell(\omega^2) = i[\Sigma_\ell^R(\omega^2) - \Sigma_\ell^A(\omega^2)].
\end{equation}

\paragraph{Convergence rate and concrete error bound.}
The decimation can be written as a transfer-matrix recursion. Eliminating
odd-numbered layers in the bulk lead leads to an effective
coupling between layers separated by $2,4,8,\ldots$ unit cells. In
matrix form, writing $T_n$ for the effective transfer operator after
$n$ decimation steps,
\begin{equation}\label{eq:Tn_recursion}
  T_{n+1} = T_n^2,
  \qquad
  T_0 = -\epsilon_0^{-1} H_{01}.
\end{equation}
The eigenvalues of $T_0$ are $\lambda_j = e^{i k_j a}$ with $k_j$ the
complex band wavevectors at energy $\omega^2$. In a bulk gap,
$|\lambda_j| < 1$ strictly for evanescent solutions and
$|\lambda_j|\to\infty$ for those that grow into the lead; only the
former survive the iteration. After $n$ steps,
\begin{equation}\label{eq:sr_decay}
  \|\alpha_n\|,\;\|\beta_n\| \;\sim\; \rho^{2^n},
  \qquad
  \rho := \max_j\bigl\{|\lambda_j|: |\lambda_j|<1\bigr\} < 1,
\end{equation}
i.e.  double-exponential convergence in $n$. On a
propagating band, $\max_j |\lambda_j| = 1$ exactly, and the convergence
is entirely due to $\eta>0$: writing
$|\lambda_j(\omega^2 + i\eta)| = 1 - c_j \eta + \mathcal{O}(\eta^2)$,
\begin{equation}\label{eq:sr_decay_band}
  \|\alpha_n\|\sim e^{-c_j\,\eta\cdot 2^n},
\end{equation}
which is still exponential in $2^n$ but with a small effective rate.
At van Hove singularities where $v_g \to 0$, $c_j \to 0$ and the
decay degenerates to algebraic.

\paragraph{Alternative algorithms.}
Several alternatives exist. Direct eigenvalue methods solve the
quadratic eigenvalue problem
$(\lambda^2 H_{10} + \lambda(H_{00} - \omega^2 \mathbb{I}) +
H_{01})\xi = 0$ to obtain $g_s$ in closed form~\cite{ando_quantum_1991,ong_efficient_2015},
and have been claimed to be faster than Sancho-Rubio for
electron transport~\cite{ong_efficient_2015}. for phonons they have the additional
advantage of producing mode-resolved transmission directly.
\todo[inline]{Elaborate on the others in quatrex}

% ----------------------------------------------------------------------
\subsection{Phonon--phonon scattering self-energy}
\label{sub:phph_se}
% ----------------------------------------------------------------------

The anharmonic self-energy has been reported in slightly different forms in the
literature: the prefactor reported for the same diagram in~\citet{wang_nonequilibrium_2007,wang_quantum_2008},
\citet{mingo_anharmonic_2006}, and \citet{luisier_atomistic_2012} differ by
factors of $\pi/9$ and $4$, as analysed by~\citet{guo_quantum_2020}
who also provided the diagrammatic derivation.
We give a derivation below.

Let $\hat{V}_3(\tau) = \tfrac{1}{3!}\sum_{\mu\nu\lambda}
\Phi_{\mu\nu\lambda}\,\tilde{u}_\mu(\tau)\,\tilde{u}_\nu(\tau)\,\tilde{u}_\lambda(\tau)$
be the cubic interaction in the interaction picture (With respect to the harmonic operators). For the interaction
we compute the S-matrix~\cite{rammer_quantum_1986,haug_quantum_2008}
\begin{equation}\label{eq:S_matrix}
  S_{\mathcal{C}} = T_{\mathcal{C}}\exp\!\left[-\frac{i}{\hbar}\oint_{\mathcal{C}}d\tau\,\hat{V}_3(\tau)\right].
\end{equation}
Wick's theorem (valid for the contour-ordered Green's function~\cite{wang_quantum_2008}) decomposes contour-ordered
products into pair contractions. The full Green's function is
\begin{equation}\label{eq:G_full}
  G_{\mu\nu}(\tau,\tau')
  = -\frac{i}{\hbar}\,\frac{
       \langle T_{\mathcal{C}}\,
       S_{\mathcal{C}}\,
       \tilde{u}_\mu(\tau)\,
       \tilde{u}_\nu(\tau')\rangle_0}
      {\langle S_{\mathcal{C}}\rangle_0},
\end{equation}
and the linked-cluster theorem cancels the denominator against the
disconnected diagrams in the numerator, leaving only connected diagrams
in the self-energy.

% ----------------------------------------------------------------------
%  (a) Sunset diagram
% ----------------------------------------------------------------------
\begin{figure}[ht]
\centering
\begin{tikzpicture}[every node/.style={font=\small}]
  \coordinate (v1) at (0,0);  \coordinate (v2) at (4,0);
  \draw[phonon] (-1.8,0) -- (v1);
  \draw[phonon] (v2)      -- (5.8,0);
  \node[above] at (-1.0, 0.05) {$\mu,\tau$};
  \node[above] at ( 4.95,0.05) {$\mu',\tau'$};
  \draw[phonon] (v1) to[bend left=55]  (v2);
  \draw[phonon] (v1) to[bend right=55] (v2);
  \node at (2, 1.65) {$G_{\nu_1\nu_4}$};
  \node at (2,-1.65) {$G_{\nu_2\nu_3}$};
  \node[fc3 vertex] at (v1) {};
  \node[fc3 vertex] at (v2) {};
  \node[below=18pt] at (v1) {$\Phi_{\mu\nu_1\nu_2}$};
  \node[below=18pt] at (v2) {$\Phi_{\mu'\nu_3\nu_4}$};
\end{tikzpicture}
\caption{Sunset (two-vertex bubble) Feynman diagram for the leading
$(\Phi^{(3)})^2$ phonon--phonon self-energy of
\cref{eq:sigma_sunset_contour}. Wavy lines denote phonon propagators
$G$, filled circles denote cubic FC3 vertices. The two equivalent
internal lines produce the symmetry factor $\tfrac{1}{2}$.}
\label{fig:sunset_diagram}
\end{figure}


To order $(\Phi^{(3)})^2$, the proper self-energy is given by the
``sunset'' diagram (\cref{fig:sunset_diagram}), in which two cubic
vertices are connected by two internal phonon propagators. The
value of the diagram, including all combinatorial factors,
is~\cite{guo_quantum_2020,wang_nonequilibrium_2014}
\begin{equation}\label{eq:sigma_sunset_contour}
  \Sigma^{(2)}_{\mu\mu'}(\tau,\tau')
  = \frac{i\hbar}{2}\sum_{\nu_1\nu_2\nu_3\nu_4}
    \Phi_{\mu\nu_1\nu_2}\,
    G_{\nu_1\nu_4}(\tau,\tau')\,
    G_{\nu_2\nu_3}(\tau,\tau')\,
    \Phi_{\mu'\nu_3\nu_4}.
\end{equation}

\begin{figure}[ht]
\centering
\begin{tikzpicture}[every node/.style={font=\small}]
  \coordinate (v) at (0,0);
  \draw[phonon] (-2.5,0) -- (v);
  \draw[phonon] (v)      -- (2.5,0);
  \node[above] at (-2.0, 0.05) {$\mu,\tau$};
  \node[above] at ( 2.0, 0.05) {$\mu',\tau'$};
  \draw[phonon] (v) to[bend left=80, looseness=2.0] (0,2.4);
  \draw[phonon] (0,2.4) to[bend left=80, looseness=2.0] (v);
  \node at (0,2.95) {$G^{<}_{\nu\lambda}$};
  \node[fc4 vertex] at (v) {};
  \node[below=12pt] at (v) {$\Phi^{(4)}_{\mu\mu'\nu\lambda}$};
\end{tikzpicture}
\caption{Hartree--Fock (one-loop) tadpole self-energy from a single
quartic FC4 vertex. The filled square denotes
$\Phi^{(4)}_{\mu\mu'\nu\lambda}$. The two internal lines contract into
the closed loop $G^{<}_{\nu\lambda}(\tau,\tau)$. The diagram is real
and frequency-independent and acts as a temperature-dependent
renormalisation $D\mapsto D + \Sigma^{(4),\mathrm{HF}}$.}
\label{fig:fc4_tadpole}
\end{figure}


A second diagram, lower order in derivatives but one loop, survives:
the FC4 Hartree--Fock ``loop'' tadpole
$\Sigma^{(4),\mathrm{HF}}_{\mu\mu'} = -\tfrac{\hbar}{2}
\sum_{\nu\lambda} \Phi^{(4)}_{\mu\mu'\nu\lambda}\,
G^{<}_{\nu\lambda}(\tau,\tau)$ from a single quartic vertex (\cref{fig:fc4_tadpole}).
The loop self-energy is real and frequency-independent.
It therefore produces no phonon linewidth, but shifts the phonon
dispersion, $D \mapsto D + \Sigma^{(4),\mathrm{HF}}$, providing the
leading temperature-dependent mode renormalisation of self-consistent
phonon theory~\cite{tadano_first-principles_2018,wang_nonequilibrium_2014}. In a
transport-focused SCBA targeting the imaginary part of $\Sigma_s$, the
loop can be absorbed into a renormalised harmonic spectrum
computed independently (e.g.\ via SCP or
SSCHA~\cite{paulatto_first-principles_2015,tadano_first-principles_2018}) and is not
re-included diagram-by-diagram in the iteration.

The factor $1/2$ above corresponds to the fully-symmetric FC3
$\Phi_{\mu\nu\lambda}$ defined by
$\Phi_{\mu\nu\lambda} = \partial^3 V/(\partial u_\mu \partial u_\nu \partial u_\lambda)$,
divided by $3!$ in \cref{eq:H_phonon}. ~\citet{mingo_anharmonic_2006} obtains an extra factor of $\pi/9$
which~\citet{guo_quantum_2020} attribute to a different convention in
the symmetrization of the cubic vertex.

Applying the Langreth analytic continuation rules~\cite{devreese_linear_1976,haug_quantum_2008}
to the contour product of two $G$'s in
\cref{eq:sigma_sunset_contour} yields real-time lesser/greater
components,
\begin{equation}\label{eq:sigma_guo}
  \Sigma^{\gtrless}_{\mu\mu'}(\omega)
  = \frac{i\hbar}{2}
    \sum_{\nu_1\nu_2\nu_3\nu_4}
    \Phi_{\mu\nu_1\nu_2}
    \left[
      \int\frac{d\omega'}{2\pi}\,
      G^{\gtrless}_{\nu_1\nu_4}(\omega')\,
      G^{\gtrless}_{\nu_2\nu_3}(\omega-\omega')
    \right]
    \Phi_{\mu'\nu_3\nu_4},
\end{equation}
which is the form derived in detail in~\cite{wang_nonequilibrium_2014,guo_quantum_2020}.
The convolution enforces the
energy-conservation $\delta(\omega - \omega' - \omega'')$ implicit in
the contour integral. Substituting equilibrium FDT
\cref{eq:phonon_fdt} reproduces Fermi's-golden-rule expressions for
three-phonon scattering rates: in mode basis,
\begin{align}\label{eq:fgr}
  \mathrm{Im}\,\Sigma^R_s(\omega_s,\qp)
  &\;=\;
   -\,\frac{\pi\hbar}{8 \omega_s(\qp)}
   \sum_{s_1 s_2}\sum_{\qp{}'}
   |V_3(s\qp; s_1\qp{}'; s_2,\qp{-}\qp{}')|^2 \notag\\
  &\quad\times
   \Bigl\{
     [1 + n_1 + n_2]\,\delta(\omega_s - \omega_{s_1} - \omega_{s_2})
     + 2[n_1 - n_2]\,\delta(\omega_s - \omega_{s_1} + \omega_{s_2})
   \Bigr\},
\end{align}
with $n_i = n_{\mathrm{BE}}(\omega_{s_i},T)$, exhibiting the
splitting/recombination ($1\to 2$) channel and the absorption ($2\to 1$)
channel. This expression coincides with the Boltzmann result of~\citet{klemens_anharmonic_1966}
in the relaxation-time approximation,.

\subsubsection{Retarded self-energy}
\label{subsub:KK}

By construction $\Sigma^R$ is causal in time, so its frequency
components satisfy a Kramers--Kronig relation. Starting from
$\Sigma^R = \tfrac{1}{2}(\Sigma^t + \Sigma^{\bar{t}}) +
\tfrac{1}{2}(\Sigma^> - \Sigma^<)$ and using the spectral identity for
$\Sigma^{t,\bar{t}}$, one obtains~\cite{guo_quantum_2020,wang_nonequilibrium_2014}
\begin{equation}\label{eq:sigma_R_full}
  \Sigma^R(\omega)
  = \tfrac{1}{2}\bigl[\Sigma^>(\omega) - \Sigma^<(\omega)\bigr]
    + i\,\mathcal{P}\!\!\int_{-\infty}^{\infty}\!\frac{d\omega'}{2\pi}\,
      \frac{\Sigma^>(\omega') - \Sigma^<(\omega')}{\omega - \omega'}.
\end{equation}
The antihermitian (dissipative) first term carries inverse-lifetime
information,
\begin{equation}\label{eq:lifetime}
  \tau_s^{-1}(\qp)
  \equiv -\frac{2}{\hbar\omega_s(\qp)}\,
     \mathrm{Im}\,\langle s|\Sigma^R(\omega_s,\qp)|s\rangle,
\end{equation}
while the dispersive second term encodes the anharmonic frequency
renormalization $\Delta\omega_s \approx
\mathrm{Re}\,\Sigma^R/(2\omega_s)$ (mode hardening or softening).
Numerical evaluation of the principal value is expensive, so in practice
the dispersive part is often dropped:
\begin{equation}\label{eq:sigma_r}
  \Sigma^R(\omega)
  \approx \tfrac{1}{2}\bigl[\Sigma^>(\omega) - \Sigma^<(\omega)\bigr].
\end{equation}

\subsubsection{Detailed balance}

The $\Sigma^{\gtrless}$ of \cref{eq:sigma_guo} satisfy detailed balance
at equilibrium,
\begin{equation}\label{eq:detailed_balance}
  \Sigma^>(\omega) = e^{\beta\hbar\omega}\,\Sigma^<(\omega),
\end{equation}
which together with the Keldysh equation \cref{eq:keldysh_eq}
guarantees that $G^{\lessgtr}$ also satisfy detailed balance, recovering
\cref{eq:phonon_fdt} when $\Sigma_L,\Sigma_R$ correspond to a single
temperature. The presence of two leads at different temperatures breaks
detailed balance and drives a steady-state heat current.

\subsubsection{Higher-order corrections}
\label{subsub:higher_order}
\begin{figure}[ht]
\centering
\begin{tikzpicture}[every node/.style={font=\small}]
  \coordinate (v1) at (0,0);     \coordinate (v2) at (5,0);
  \coordinate (vt) at (2.5, 1.4);  \coordinate (vb) at (2.5,-1.4);
  \draw[phonon] (-1.8,0) -- (v1);
  \draw[phonon] (v2)      -- (6.8,0);
  \node[above] at (-1.0, 0.05) {$\mu,\tau$};
  \node[above] at ( 5.95,0.05) {$\mu',\tau'$};
  \draw[phonon] (v1) -- (vt);   \draw[phonon] (v1) -- (vb);
  \draw[phonon] (vt) -- (v2);   \draw[phonon] (vb) -- (v2);
  \draw[phonon] (vt) -- (vb);
  \foreach \pt in {v1,v2,vt,vb}{\node[fc3 vertex] at (\pt) {};}
\end{tikzpicture}
\caption{Two-loop $\Phi^{(3)}$ vertex correction to the phonon
self-energy ($\mathcal{O}((\Phi^{(3)})^4)$). Four cubic vertices form a
1PI ``diamond'' subgraph}
\label{fig:vertex_correction}
\end{figure}

\begin{figure}[ht]
\centering
\begin{subfigure}[b]{0.95\linewidth}
\centering
\begin{tikzpicture}[every node/.style={font=\small}]
  \coordinate (v1) at (0  ,0);  \coordinate (v2) at (1.6,0);
  \coordinate (v3) at (3.2,0);  \coordinate (v4) at (4.8,0);
  \draw[phonon] (-1.8,0) -- (v1);
  \draw[phonon] (v4)      -- (6.6,0);
  \node[above] at (-1.0,0.05) {$\mu,\tau$};
  \node[above] at ( 5.7,0.05) {$\mu',\tau'$};
  \draw[phonon] (v1) -- (v2);
  \draw[phonon] (v2) -- (v3);
  \draw[phonon] (v3) -- (v4);
  \draw[phonon] (v1) to[bend left=70, looseness=1.1] (v3);
  \draw[phonon] (v2) to[bend left=70, looseness=1.1] (v4);
  \foreach \pt in {v1,v2,v3,v4}{\node[fc3 vertex] at (\pt) {};}
\end{tikzpicture}
\caption{Crossed two-loop self-energy graph at order
$(\Phi^{(3)})^4$.}
\label{fig:crossed_graph}
\end{subfigure}

\vspace{1em}
\begin{subfigure}[b]{0.95\linewidth}
\centering
\begin{tikzpicture}[every node/.style={font=\small}]
  \coordinate (v1) at (0  ,0);   \coordinate (v2) at (1.8,0);
  \coordinate (v3) at (3.4,0);   \coordinate (v4) at (5.2,0);
  \draw[phonon] (-1.8,0) -- (v1);
  \draw[phonon] (v4)      -- (7.0,0);
  \node[above] at (-1.0,0.05) {$\mu,\tau$};
  \node[above] at ( 6.1,0.05) {$\mu',\tau'$};
  \draw[phonon] (v1) to[bend left=80]  (v2);
  \draw[phonon] (v1) to[bend right=80] (v2);
  \draw[phonon] (v3) to[bend left=80]  (v4);
  \draw[phonon] (v3) to[bend right=80] (v4);
  \draw[phonon] (v2) -- (v3);
  \foreach \pt in {v1,v2,v3,v4}{\node[fc3 vertex] at (\pt) {};}
\end{tikzpicture}
\caption{Ladder topology at order $(\Phi^{(3)})^4$:
two sunset bubbles connected by a single line.}
\label{fig:ladder_graph}
\end{subfigure}

\caption{Non-trivial two-loop self-energy topologies at order
$(\Phi^{(3)})^4$ that lie beyond bubble-SCBA.}
\label{fig:two_loop_topologies}
\end{figure}

\Cref{eq:sigma_guo} captures the leading three-phonon process. Higher
orders include~\cite{wang_nonequilibrium_2014,guo_quantum_2020}:
(i) the four-phonon tadpole from $\Phi^{(4)}$,
$\Sigma^{(4),\mathrm{HF}}_{\mu\mu'} = -\tfrac{\hbar}{2}\sum_{\nu\lambda}
\Phi^{(4)}_{\mu\mu'\nu\lambda}\,G^<_{\nu\lambda}(\tau,\tau)$,
which is real, frequency-independent, and renormalizes
$D \mapsto D + \Sigma^{(4),\mathrm{HF}}$, contributing the
temperature-dependent thermal expansion and frequency
shifts~\cite{wang_nonequilibrium_2014} (\cref{fig:fc4_tadpole}),
(ii) the two-loop $\Phi^{(3)}$ vertex correction, relevant at high
$T$ where $\hbar\omega_{\max}/k_B T \lesssim 1$ and three-phonon
processes saturate (\cref{fig:vertex_correction}) and
(iii) crossed and ladder graphs that are needed to maintain
$\Phi$-derivability.

The $\Phi$-derivability framework of Baym and Kadanoff~\cite{baym_conservation_1961, baym_self-consistent_1962}
states that $\Sigma$ must be the functional derivative of a generating
functional $\Phi[G]$ for the conservation laws (energy, particle number)
to be satisfied identically. The sunset diagram of
\cref{eq:sigma_sunset_contour} is $\Phi$-derivable when the same
$G$ enters both internal lines, but only asymptotically as the
SCBA iteration converges~\cite{lee_anharmonic_2018}. This is e.g. addressed by
the LOA-based scheme of~\citet{lee_anharmonic_2018}
which is exactly conserving at every iteration.

\subsubsection{Transverse periodic case}

For a 3D system with periodicity transverse to transport, the
self-energy acquires a $\qp$ argument and an internal sum over
transverse momenta~\cite{guo_quantum_2020}:
\begin{align}\label{eq:sigma_phph}
  \sse{\gtrless}_{x x'}(\omega,\qp)
  = \frac{i\hbar}{2}
    &\sum_{y_1 y_2 y_3 y_4}
    \frac{1}{N}\sum_{\qp{}'}
    \tilde{\Phi}_{x y_1 y_2}\!\bigl(\qp{}',\qp - \qp{}'\bigr)\,
    \tilde{\Phi}_{x' y_3 y_4}\!\bigl(\qp{}'-\qp,-\qp{}'\bigr)
    \notag\\
  &\times
    \int\!\dw\;
    G^{\gtrless}_{y_1 y_4}(\omega',\qp{}')\,
    G^{\gtrless}_{y_2 y_3}(\omega-\omega',\qp - \qp{}'),
\end{align}
where $x,y$ denote transport-direction DOF indices, i.e.\ $x=(l_x,\mu)=(l_x,\kappa,\alpha)$
with $l_x$ enumerating transport-direction slabs/supercells. The
partially Fourier-transformed FC3 entering this expression is
\begin{align}\label{eq:fc3_fourier}
  \tilde{\Phi}_{x y_1 y_2}(\qp,\qp{}')
  = \tilde{\Phi}^{l m_1 m_2}_{\mu \nu_1 \nu_2}(\qp,\qp{}')
  = \sum_{\substack{n_1,\,n_2\\ n_{ix}=m_{ix}}}
    \Phi^{l n_1 n_2}_{\mu \nu_1 \nu_2}\;
    e^{-i\qp\cdot (n_1 - m_1)}\;
    e^{-i\qp{}'\cdot (n_2 - m_2)},
\end{align}
i.e.\ the sum runs over transverse supercells at fixed transport-direction
cell. Unlike the harmonic FC2 object which depends on a single
transverse wavevector, the cubic tensor depends on two: the
$\qp{}'$ sum couples all transverse wavevectors, so evaluating
$\Sigma_s(\qp)$ requires Green's functions at every $\qp{}'$, destroying
the embarrassing parallelism of the ballistic problem. The arithmetic
complexity scales as $\mathcal{O}(N_q^2 N_\omega^2 N_{\mathrm{DOF}}^4)$
before any low-rank or stochastic acceleration. This is the dominant
cost of the anharmonic simulation, as reported in~\cite{luisier_atomistic_2012,guo_quantum_2020, lee_anharmonic_2018}.

% ----------------------------------------------------------------------
\subsection{Observables}
\label{sub:negf_obs}
% ----------------------------------------------------------------------

We collect the observables computable from $G^{R,A,\lessgtr}$ in both
the ballistic and anharmonic regimes.

\subsubsection{Spectral function and local density of states}

The spectral function and LDOS are
\begin{equation}\label{eq:dos}
  A(\omega^2) = i[G^R - G^A],
  \qquad
  \rho_\mu(\omega^2) = \frac{1}{\pi}\,A_{\mu\mu}(\omega^2),
\end{equation}
related to the standard $\omega$-resolved density of states by
$\rho(\omega) = 2\omega\,\rho(\omega^2)$. In the transverse periodic
case $\rho_\mu(\omega) = (2\omega/N_q)\sum_{\qp} A_{\mu\mu}(\omega^2,\qp)/\pi$.

\subsubsection{Ballistic transmission and conductance}

In the ballistic limit ($\Sigma_s = 0$) the transmission function is
given by the Caroli formula~\cite{caroli_direct_1971,mingo_calculation_2003,wang_nonequilibrium_2007},
\begin{equation}\label{eq:caroli}
  \mathcal{T}(\omega^2,\qp)
  = \mathrm{Tr}\bigl[\Gamma_L\,G^R\,\Gamma_R\,G^A\bigr],
\end{equation}
with the transverse-averaged transmission
\begin{equation}
  \bar{\mathcal{T}}(\omega^2)
  = \frac{1}{N_q}\sum_{\qp} \mathcal{T}(\omega^2,\qp).
\end{equation}
The thermal conductance per unit area is the Landauer integral
\begin{equation}\label{eq:G_thermal}
  G(T)
  = \frac{1}{A_c}\int_0^\infty\!\frac{d\omega}{2\pi}\;
    \bar{\mathcal{T}}(\omega)\,\hbar\omega\,
    \frac{\partial n_\mathrm{BE}}{\partial T},
\end{equation}
with $A_c$ the cross-sectional area and
\begin{equation}
  \frac{\partial n_\mathrm{BE}}{\partial T}
  = \frac{\hbar\omega}{k_B T^2}\,
    \frac{e^{\hbar\omega/k_B T}}
         {(e^{\hbar\omega/k_B T}-1)^2}.
\end{equation}
The integrand vanishes at both endpoints, so a trapezoidal rule on the
sampled frequency grid suffices.

\subsubsection{Mode-resolved transmission}

Diagonalizing $\Gamma_\ell$ at fixed $(\omega,\qp)$ yields incoming and
outgoing mode bases of the leads. Writing $\Gamma_L = U_L \Lambda_L
U_L^\dagger$ and similarly for $R$, the transmission matrix element from
incoming mode $m\in L$ to outgoing mode $m'\in R$ is
\begin{equation}\label{eq:mode_T}
  t_{m'm}(\omega,\qp)
  = \sqrt{\Lambda_R^{m'}}\,
    \bigl[U_R^\dagger G^R U_L\bigr]_{m'm}\,
    \sqrt{\Lambda_L^{m}},
\end{equation}
with $\sum_{m'm}|t_{m'm}|^2 = \mathcal{T}$. This is the phonon analogue
of the Fisher--Lee relation~\cite{fisher_relation_1981}.
For phonons the mode-matching formulation~\cite{ong_efficient_2015}
provides an alternative route to the same
quantity through the eigenproblem of the bulk transfer matrix.

\subsubsection{Meir--Wingreen heat current}

In the anharmonic regime the Caroli formula no longer holds: there is
no $\mathcal{T}(\omega)$ that captures both elastic and inelastic
processes. The generalization is the phonon Meir--Wingreen
formula~\cite{meir_landauer_1992,wang_nonequilibrium_2007,wang_quantum_2008},
\begin{equation}\label{eq:meir_wingreen}
  I_L
  = \int_0^\infty\!\frac{d\omega}{4\pi}\,\hbar\omega\,
    \mathrm{Tr}\bigl[
      \Sigma_L^<(\omega)\,G^>(\omega)
    - \Sigma_L^>(\omega)\,G^<(\omega)
    \bigr].
\end{equation}
\Cref{eq:meir_wingreen} is exact for any
many-body $\Sigma_s$, provided $G^{\lessgtr}$ are obtained from the full
Keldysh equation \cref{eq:keldysh_eq}. In the ballistic limit, the
Keldysh equation gives $G^{<} = G^R(\Sigma_L^< + \Sigma_R^<)G^A$, and
substitution into \cref{eq:meir_wingreen} reproduces \cref{eq:caroli}
weighted by $\hbar\omega[n_L-n_R]$, recovering Landauer.

\subsubsection{Local heat current and conservation}

The bond-resolved heat current of Hardy~\cite{hardy_energy-flux_1963} can provide atomic insight.
Between two transport-direction slabs $i$
and $j = i + 1$,
\begin{equation}\label{eq:I_local}
  I_{i,i+1}
  = -\int_0^\infty\!\frac{d\omega}{2\pi}\,\hbar\omega\,
    \mathrm{Re}\,\mathrm{Tr}\bigl[H_{i,i+1}\,G^<_{i+1,i}(\omega)\bigr],
\end{equation}
with $H_{i,i+1}$ the inter-slab coupling block. In a converged SCBA
solution the steady-state continuity equation enforces $I_{i,i+1}$
independent of $i$ throughout the device, providing physically meaningful convergence test:
\begin{equation}
  \delta I = \max_i I_{i,i+1} - \min_i I_{i,i+1}
\end{equation}
measures the violation of energy conservation by approximations in
$\Sigma_s$ (incomplete SCBA iteration, Hilbert-transform truncation,
finite frequency grid). A convergence criteron for SCBA can be
$\delta I/\bar I \lesssim \varepsilon$~\cite{luisier_atomistic_2012,guo_quantum_2020} for
appropriate $\varepsilon$ (e.g. $10^{-3}$ in \cite{guo_quantum_2020}).

Heat conservation at the level of \cref{eq:sigma_sunset_contour}
is guaranteed by $\Phi$-derivability~\cite{baym_self-consistent_1962} only at full
self-consistency. At any finite iteration count, $\delta I/\bar I$ is
non-zero. The Hilbert-transform neglect of \cref{eq:sigma_r} introduces an additional
$\mathcal{O}(\Delta\omega/\omega)$ conservation error.

% ----------------------------------------------------------------------
\subsection{SCBA iterations}
\label{sub:scba}
% ----------------------------------------------------------------------

\Cref{eq:system,eq:keldysh_eq,eq:sigma_phph,eq:sigma_r} together define
a nonlinear fixed-point problem $\Sigma_s = \mathcal{F}[G[\Sigma_s]]$,
solved iteratively in the SCBA. Convergence is two-fold: numerical
(consecutive iterates close in some norm) and physical (current
conservation, \cref{eq:I_local}). Mixing schemes are necessary because
pure fixed-point iteration of \cref{eq:sigma_phph} is often divergent or
oscillatory at finite anharmonicity~\cite{wang_nonequilibrium_2007,luisier_atomistic_2012}.

\subsubsection{Linear mixing}

The simplest scheme,
\begin{equation}\label{eq:linear_mix}
  \Sigma^{(n+1)} = (1-\alpha)\,\Sigma^{(n)} + \alpha\,\Sigma^{\mathrm{new}},
\end{equation}
with damping parameter $\alpha \in (0,1]$. Linear mixing converges if
the fixed-point Jacobian
$J = \partial\mathcal{F}/\partial\Sigma_s$ has spectral radius
$\sigma(I - \alpha J) < 1$, i.e.\ $\alpha < 2/(1+|\sigma(J)|_{\max})$.
For weak anharmonicity ($|\sigma(J)|<1$) any $\alpha \le 1$ works; for
strong anharmonicity the spectrum of $J$ extends close to $-1$ in
some directions and small $\alpha\sim 0.1\text{--}0.3$ is required, with
linear convergence rate $1-\alpha$ in the slowest mode.

\subsubsection{Anderson mixing}

Anderson acceleration~\cite{anderson_iterative_1965} constructs a quasi-Newton
update from a history of residuals,
\begin{equation}\label{eq:anderson}
  \vc{x}_{k+1}
  = \bigl(\vc{x}_k + \alpha\,\vc{f}_k\bigr)
    - \bigl(\Delta\vc{X} + \alpha\,\Delta\vc{F}\bigr)\,\vc{\gamma}_k,
\end{equation}
with $\vc{f}_k = G(\vc{x}_k) - \vc{x}_k$ the residual,
$\Delta\vc{X}, \Delta\vc{F}$ the matrices of consecutive iterate and
residual differences, and $\vc{\gamma}_k$ solving the regularized
least-squares problem
\begin{equation}
  (\Delta\vc{F}^H \Delta\vc{F} + \lambda \mathbb{I})\,\vc{\gamma}
  = \Delta\vc{F}^H \vc{f}_k,
  \qquad
  \lambda = 10^{-8}\,\tr(\Delta\vc{F}^H \Delta\vc{F}).
\end{equation}
For linear problems Anderson mixing is
essentially equivalent to GMRES applied to the residual
equation~\cite{walker_anderson_2011,rohwedder_analysis_2011}, with provable
$q$-linear convergence with $q$-factor $q = 3c - c^2$ for depth-1
Anderson when the underlying iteration is a contraction with constant
$c < 2-\sqrt{3}$~\cite{toth_convergence_2015,bian_anderson_2021}. Local
convergence in nonlinear electronic-structure problems is well
established under mild Lipschitz assumptions on the
fixed-point map~\cite{rohwedder_analysis_2011,toth_convergence_2015}. In
practice, Anderson mixing dramatically accelerates convergence when the
fixed-point equation is well-conditioned, with rates approaching that
of Newton's method. However, storing the history is often prohibitive, so that
Anderson mixing becomes infeasible for large structures.

\subsubsection{Convergence diagnostics}

We monitor three quantities across iterations: (i) the relative
Frobenius norm
$\|\Sigma_s^{(n+1)}-\Sigma_s^{(n)}\|/\|\Sigma_s^{(n)}\|$,
(ii) the conductance $G(T)$ from \cref{eq:G_thermal} or, in the
inelastic case, the Meir--Wingreen current, and (iii) the slab-resolved
current imbalance $\delta I/\bar I$. Convergence is declared when (i)
and (iii) fall below prescribed tolerances simultaneously.

\subsubsection{Lowest-order approximation.}

An attractive alternative to direct SCBA, developed by~\citet{lee_anharmonic_2018,lee_quantum_2019}, is the
lowest-order approximation (LOA) implemented through Padé
resummation. LOA collects only the conserving ($\Phi$-derivable) diagrams
at each order in $\Phi^{(3)}$, ensuring exact current conservation
at every iteration, then accelerates the resulting series by
Padé approximants. Compared to SCBA, LOA achieves comparable accuracy
at $\sim 2\times$ lower cost for thermal transport in 3D
nanostructures~\cite{lee_anharmonic_2018}.

% ----------------------------------------------------------------------
\subsection{Decay properties and Truncation errors}
\label{sub:decay}
% ----------------------------------------------------------------------

Truncations at every step (real-space FC cutoff, finite device length, finite $\omega$ grid,
finite Sancho-Rubio iteration count) introduce errors with
quantifiable scaling. We summarize, from first principles, the decay
rates of each of the involved quantities.

\subsubsection{Real-space decay of harmonic force constants}
\label{subsub:fc2_decay}

The asymptotic behaviour of $\Phi^{(2)}(R)$ at large $|R|$ is controlled
by the analytic structure of the dielectric response of the underlying
electronic system. \Citet{cances_decay_2025} prove the following
within the reduced Hartree--Fock (or random-phase
approximation) model:

\paragraph{Insulators at zero electronic temperature.}
The dynamical matrix $D(\qp)$ is singular at $\qp = 0$, with the
non-analytic dipole--dipole structure given in~\cite{cochran_dielectric_1962}:
\begin{equation}\label{eq:fc2_dipole}
  D^{\mathrm{LR}}_{\kappa\alpha,\kappa'\beta}(\qp \to 0)
  \;\sim\;
  \frac{4\pi}{\Omega_{\mathrm{cell}}}\,
  \frac{(\qp\cdot Z^*_{\kappa})_\alpha\,(\qp\cdot Z^*_{\kappa'})_\beta}
       {\qp\cdot \epsilon_\infty\cdot \qp},
\end{equation}
where $Z^*_\kappa$ is the Born effective charge tensor of atom $\kappa$
and $\epsilon_\infty$ the optical dielectric tensor. Inverse Fourier
transform of this $\qp\to 0$ singularity gives an algebraic
real-space tail to $\Phi$.

\paragraph{Materials at finite electronic temperature (metals).}
\Citet{cances_decay_2025} prove that the
dynamical matrix is analytic in $\qp$, and
$\Phi^{(2)}(R)$ decays exponentially:
\begin{equation}\label{eq:fc2_metal}
  |\Phi^{(2)}(R)|\;\le\;C\,e^{-|R|/\xi_e},\qquad
  \xi_e \;\sim\;\frac{\hbar v_F}{k_B T_{\mathrm{el}}}.
\end{equation}

\paragraph{Concrete truncation error.}
Truncating at distance $R_c$ introduces an error in the dynamical matrix
of order
\begin{equation}\label{eq:fc2_error_metal}
  \|D(\qp) - D^{R_c}(\qp)\| \;\lesssim\;C \frac{e^{-R_c/\xi_e}}{1 - e^{-1/\xi_e}}
  \quad\text{(metals, exponential)},
\end{equation}
or, in polar insulators, $\sim 1/R_c^3$ for the bare truncation, but
$\sim e^{-R_c/\xi_e}$ if the short-range part once the
$1/r^3$ Cochran--Cowley tail has been subtracted analytically and added
back exactly via Ewald summation~\cite{gonze_dynamical_1997}.

\subsubsection{Transverse $\qp$-mesh convergence}
\label{subsub:qp_decay}

The convergence of observables in $N_q$ depends on the differentiability
of the integrand. For the Caroli formula \cref{eq:caroli}, the
$\qp$-integrand is generically smooth and trapezoidal Monkhorst--Pack
sampling converges as $\mathcal{O}(N_q^{-2})$ in 2D. At Brillouin-zone boundaries or near van
Hove singularities, convergence degrades to $\mathcal{O}(N_q^{-1})$;
a Gaussian smearing or tetrahedron method may then be preferable. The
self-energy convolution \cref{eq:sigma_phph} couples $\qp$ and $\qp{}'$,
so the effective integrand has $N_q^{-1/2}$ statistical-noise structure
under random sampling, but adaptive importance sampling can recover
$\mathcal{O}(N_q^{-1})$ ~\cite{lee_anharmonic_2018}.

\subsubsection{Convergence rate of Sancho-Rubio decimation revisited}

As discussed in \cref{sub:contact_se}, the decimation iteration
\cref{eq:sr_eps_s,eq:sr_eps,eq:sr_alpha,eq:sr_beta} effectively doubles
the number of decimated layers at each step, with iterates obeying
$\|\alpha_n\|, \|\beta_n\| \sim \rho^{2^n}$ inside gaps and
$e^{-c\eta\cdot 2^n}$ on bands. The practical implication is that
\emph{doubling the tolerance roughly costs one extra iteration}; total
cost is $\mathcal{O}(N_\omega\,N_q\,N_{\mathrm{DOF}}^3\,N_{\mathrm{iter}})$
with $N_{\mathrm{iter}}$ essentially constant (and small,
$\sim 10\text{--}30$).

% ----------------------------------------------------------------------
\section{Computation of the Phonon-Phonon Scattering Self Energy}
\label{sec:phph_computation}
% ----------------------------------------------------------------------

The phonon-phonon SSE in \cref{eq:sigma_phph} involves three computational
steps that dominate the cost of an SCBA iteration of the anharmonic
SCAB equations: a four-index ring
contraction over $y_1,y_2,y_3,y_4$, an internal sum over the
transverse wavevector $\qp'$, and a frequency convolution along
$\omega$. The frequency convolution is structurally identical to the
polarisation that appears in the GW formalism
\cite{deuschle_electron-electron_2025},
\begin{equation}\label{eq:gw_polarization}
  P_{ij}^{\gtrless}(E)
  = -i \int dE'\;
       G_{ij}^{\gtrless}(E')\,G_{ji}^{\gtrless}(E'-E),
\end{equation}
and reuses the same FFT-based machinery. The $\qp'$-sum couples
distinct transverse wavevectors through the FC3 vertices
$\tilde\Phi$ and is unique to the transverse-periodic geometry. The
inner four-index contraction has no analogue in GW: in
\cref{eq:gw_polarization}, the contraction structure of $P$ is local
in the orbital pair $(i,j)$, whereas the phonon bubble couples
$y_1,\dots,y_4$ across the supercell. In what follows we omit the
$\gtrless$ superscripts unless the distinction matters; all
derivations apply unchanged to lesser, greater, and (via the
fluctuation--dissipation construction of \cref{eq:sigma_r}) retarded
self-energies.

% ----------------------------------------------------------------------
\subsection{Frequency convolution and discretisation}
\label{ssec:phph_freq}
% ----------------------------------------------------------------------

The integrand in \cref{eq:sigma_phph} contains a single frequency
convolution $\int d\omega'\;G_1(\omega')\,G_2(\omega-\omega')$. We
turn it into a pointwise product in time through the Fourier
transform pair
\begin{equation}\label{eq:g_tau}
  g(\tau;\cdot)
  = \int\!\frac{d\omega}{2\pi}\;
    G(\omega;\cdot)\,e^{-i\omega\tau},
  \qquad
  G(\omega;\cdot)
  = \int\!d\tau\;g(\tau;\cdot)\,e^{+i\omega\tau},
\end{equation}
under which the 1D convolution theorem gives
\begin{equation}\label{eq:omega_conv}
  \int\!d\omega'\;
    G_1(\omega';\cdot)\,G_2(\omega-\omega';\cdot)
  = \int\!d\tau\;
    g_1(\tau;\cdot)\,g_2(\tau;\cdot)\,e^{+i\omega\tau}.
\end{equation}
Substituting \cref{eq:omega_conv} into \cref{eq:sigma_phph}, the
bubble takes the $\tau$-domain form
\begin{equation}\label{eq:sigma_phph_tau}
\begin{aligned}
  \Sigma_{xx'}(\omega;\qp)
  = \frac{i\hbar}{2}
    &\sum_{y_1y_2y_3y_4}
    \frac{1}{N}\sum_{\qp'}
    \tilde\Phi_{xy_1y_2}(\qp',\qp-\qp')\,
    \tilde\Phi_{x'y_3y_4}(\qp'-\qp,-\qp')\\
    &\times\int\!d\tau\;
    g_{y_1y_4}(\tau;\qp')\,g_{y_2y_3}(\tau;\qp-\qp')\,
    e^{+i\omega\tau}.
\end{aligned}
\end{equation}
The $\omega$-integral has been replaced by a $\tau$-integral, the
two propagators are now multiplied pointwise, and a
final $\tau\to\omega$ Fourier transform delivers
$\Sigma$ on the original frequency grid.

\paragraph{Discrete frequency grid.} Numerically, the SSE is
evaluated on a uniform grid
$\omega_n = \omega_\mathrm{min} + n\Delta\omega$ for
$n=0,\ldots,N_\omega-1$, with broadening
$\eta = \eta_\mathrm{factor}(\Delta\omega)^{2}$ to suppress
spurious poles at the resolution scale of the grid. The convolution
is evaluated by FFT on a zero-padded grid of length
\begin{equation}\label{eq:nfft_pad}
  N_\mathrm{fft} = 2N_\omega - 1,
\end{equation}
which is the smallest length that avoids cyclic wrap-around between
the two factors when both have support of width $N_\omega$. Stages~1 and~3
(\cref{ssec:phph_distribution}) operate on $N_\mathrm{fft}$-long
arrays, the SSE on the original grid is recovered by truncation to
the first $N_\omega$ points after the inverse FFT.

\paragraph{Convolution support: $f_{\max}\!\ge\!2\omega_{\max}/2\pi$.}
The harmonic Green's functions $G^{\gtrless}(\omega)$ have spectral
support bounded by the largest phonon frequency $\omega_{\max}$ of the
device, i.e.\ the spectrum of
$H_{00}+H_{01}+H_{01}^{\dagger}$. The product
$g_1(\tau)\,g_2(\tau)$ in \cref{eq:omega_conv} has Fourier support
that is the convolution sum of the individual supports, so
$\Sigma^{\gtrless}(\omega)$ is nonzero on
\begin{equation}\label{eq:conv_support}
  \omega\in[-2\omega_{\max},\,2\omega_{\max}]
\end{equation}
and only there. The truncated linear convolution \cref{eq:nfft_pad} is
unbiased on that interval \emph{only when the grid covers it}, i.e.\
when $f_{\max}\ge 2\omega_{\max}/2\pi$. If the grid is shorter,
$f_{\max}<2\omega_{\max}/2\pi$, the bubble contributions at
$|\omega|\in(f_{\max},2\omega_{\max})$ are missing from the inverse
FFT and, on a discrete grid, the residual aliases back into the
visible window. The bubble then carries large spurious contributions
that fake a strong-coupling self-energy and drive the SCBA loop
unstable. The required grid size is therefore set by the stiffest
mode of the device, not by the thermally populated bands; in
particular, an H-terminated wire with $\omega_{\max}/2\pi\sim 64$\,THz
from the Si--H stretch needs $f_{\max}\gtrsim 130$\,THz, not the
$\sim 16$\,THz of a bulk-Si default. \Cref{sec:phonon_solver_pkg}
documents the runtime safety net that enforces this condition.

\paragraph{Negative-frequency handling.} The bosonic phonon Green's
functions live on a real frequency grid centred on $\omega=0$. Both
$\Sigma^{<}$ and $\Sigma^{>}$ involve products of $G^{\gtrless}$ at
shifted arguments. Depending on the relative sign of $\omega$ and
$\omega'$, these give two distinct convolution patterns. Following
the convention used in the GW SSE
\cite{deuschle_electron-electron_2025}, we can split each product into a
forward convolution on the positive-frequency support and a
cross-correlation against the conjugate on the negative-frequency
support. Concretely, defining
\begin{equation}
  \mathrm{conv}(f,g)[\omega]
    = \int\!d\omega'\,f(\omega')\,g(\omega-\omega'),
  \qquad
  \mathrm{corr}(f,g)[\omega]
    = \int\!d\omega'\,f(\omega')\,g(\omega'-\omega),
\end{equation}
the bubble splits as
\begin{equation}
  \Sigma^{\gtrless}(\omega) \;\propto\;
    \mathrm{conv}\bigl(g^{\gtrless},g^{\gtrless}\bigr)[\omega]
    + \mathrm{corr}\bigl(g^{\gtrless},(g^{\lessgtr})^{*}\bigr)[\omega],
\end{equation}
the second term capturing contributions where one factor's frequency
exceeds $\omega$ in absolute value. Both pieces are evaluated by FFT
on the padded grid \cref{eq:nfft_pad}. The only book-keeping is the
sign convention on the frequency axis. The retarded SSE is
constructed from $\Sigma^{>}-\Sigma^{<}$ via \cref{eq:sigma_r},
optionally with the Hilbert-transform real-part correction obtained
by another FFT-based pass.

% ----------------------------------------------------------------------
\subsection{Inner ring contraction}
\label{ssec:phph_ring}
% ----------------------------------------------------------------------

At each working point $(\tau,\qp,\qp')$, the SSE entry $\Sigma_{xx'}$
contains the four-tensor ring contraction
\begin{equation}\label{eq:inner_contraction}
  \mathcal{B}_{xx'}
  = \sum_{y_1y_2y_3y_4}
    \Phi_{xy_1y_2}\,g_{y_1y_4}\,g_{y_2y_3}\,\Phi_{x'y_3y_4},
\end{equation}
where $\Phi$ stands for either $\tilde\Phi$ at the working
$(\qp,\qp')$ or the real-space FC3, and the indices each range over
$\ndof = N_{AO} = \NB\NBS$ degrees of freedom. The graph defined by
the index structure is a four-cycle: $x \to y_1 \to y_4 \to x'$ via
$\Phi_L$ and $g$, and $x' \to y_3 \to y_2 \to x$ via $\Phi_R$ and the
second $g$. The two $\Phi$ legs and two $g$ legs alternate around
the ring.

A ring of four tensors admits
three structurally distinct evaluation orders, distinguished by
which pair of adjacent tensors is contracted first.

\paragraph{Outside-in.} Contract one $\Phi$ leg at a time, starting
from one corner of the ring. The intermediate $A_{x,c,d} = \sum_{e}
\Phi_{x,c,e}\,g_{e,d}$ has three free indices and per-entry cost
$\cO(\ndof)$. Building $A$ for all $(x,c,d)$ costs $\cO(\ndof^{4})$
and uses $\cO(\ndof^{3})$ memory; iterating around the ring gives
two more steps of the same shape, for a total cost
$\cO(\ndof^{4})$.

\paragraph{Balanced tree.} Compute the two halves
\begin{equation}\label{eq:balanced_tree}
  T_1(x)_{y_2y_4}
  = \sum_{y_1}\Phi_{xy_1y_2}\,g_{y_1y_4},
  \qquad
  T_2(x')_{y_2y_4}
  = \sum_{y_3}g_{y_2y_3}\,\Phi_{x'y_3y_4},
\end{equation}
in parallel, then contract along the shared $(y_2,y_4)$
indices,
\begin{equation}\label{eq:balanced_pair}
  \mathcal{B}_{xx'}
  = \sum_{y_2y_4}
    T_1(x)_{y_2y_4}\,T_2(x')_{y_2y_4}.
\end{equation}
Each half-tree has the same arithmetic shape as one outside-in step
($\cO(\ndof^{4})$ to build, $\cO(\ndof^{3})$ memory per piece). The
final pairing has cost $\cO(\ndof^{4})$. The total is again
$\cO(\ndof^{4})$ but now exposes two independent halves that can be
distributed across disjoint resources.

\paragraph{Inside-out.} Contract the two propagators first into the
four-index intermediate $C_{y_1y_4y_2y_3}=g_{y_1y_4}g_{y_2y_3}$, then
sandwich with $\Phi_L,\Phi_R$. This intermediate has four free
spatial indices and $\cO(\ndof^{4})$ entries; building it costs
$\cO(\ndof^{4})$ and the subsequent contraction with each $\Phi$
costs $\cO(\ndof^{5})$. The intermediate also defeats sparsity
exploitation, because the propagator pair is generically denser
than either $g$ alone. We discard this order in what follows.

Outside-in (i) and
balanced-tree (ii) have identical asymptotic cost
\begin{equation}\label{eq:cic_cost}
  \mathcal{C}_\mathrm{ic}^\mathrm{dense}
  = \cO\!\bigl(\ndof^{4}\bigr),
\end{equation}
which we take as the unit cost throughout. Balanced-tree is
preferable when the two halves can be evaluated on different ranks,
which is the case in our distribution scheme; otherwise the two are
interchangeable. Outside-in is the natural form for a single-rank
evaluation and is the form we use to express the per-block kernel
in \cref{ssec:phph_algebra}.

% ----------------------------------------------------------------------
\subsection{Q-space and real-space formulations}
\label{ssec:phph_qspace}
% ----------------------------------------------------------------------

We present two evaluation strategies of \cref{eq:sigma_phph_tau} that differ in how
they treat the $\qp'$-sum: a direct evaluation in the Brillouin zone
(``q-space'' below), or a reformulation in real space via a second
application of the convolution theorem.

\paragraph{Q-space formulation.} The direct evaluation retains the
explicit $\qp'$-loop. At each working point $(\tau,\qp,\qp')$, the
partially Fourier-transformed FC3 $\tilde\Phi_{xy_1y_2}(\qp',\qp-\qp')$
is needed. We assemble it on the fly by streaming through the
$N_\perp^{2}$ non-zero real-space FC3 entries with the corresponding
phase factors. Because $N_\perp^{2} \ll \ndof^{2}$ for realistic
FC3 cut-offs, this phase assembly is negligible compared with the
subsequent ring contraction. We never materialise the dense
$\tilde\Phi$ at every $(\qp,\qp')$. Defining the $\tau$-bubble at
fixed transverse $\qp$,
\begin{equation}\label{eq:bubble_q}
\begin{aligned}
  \mathcal{B}_{xx'}(\tau;\qp)
  \equiv \frac{1}{N}\sum_{\qp'}
    \sum_{y_1y_2y_3y_4}
    \tilde\Phi_{xy_1y_2}&(\qp',\qp-\qp')\,
    g_{y_1y_4}(\tau;\qp')\,\\
    &\times g_{y_2y_3}(\tau;\qp-\qp')\,
    \tilde\Phi_{x'y_3y_4}(\qp'-\qp,-\qp'),
\end{aligned}
\end{equation}
the SSE follows from a single $\tau\to\omega$ Fourier transform.
The total cost is
\begin{equation}\label{eq:cost_q}
  \mathcal{W}_q = N_\tau\,N_q^{\,2}\,\mathcal{C}_\mathrm{ic},
\end{equation}
with $N_\tau = N_\mathrm{fft}$ the number of frequency grid points,
$N_q$ the number of transverse Brillouin-zone samples (e.g.\
$N_q=400$ for a $20\times 20$ grid). The factor $N_q^{2}$
is the direct consequence of the two transverse momentum arguments
of $\tilde\Phi$ and is the primary cost increase of the
transverse-periodic problem over the purely local one.

\paragraph{Real-space formulation.} The $N_q^{2}$ factor in
\cref{eq:cost_q} can in principle be eliminated by applying the
convolution theorem to the $\qp'$-sum and evaluating the resulting
expression on a finite real-space lattice of $\vR$ values.

Substituting the Fourier expansion \cref{eq:fc3_fourier} of both
FC3 factors into \cref{eq:sigma_phph}:
\begin{align}
  \tilde\Phi_{xy_1y_2}(\qp',\qp-\qp')
  &= \sum_{\vc{\rho}_1,\vc{\rho}_2}
     \Phi_{xy_1y_2}(\vc{\rho}_1,\vc{\rho}_2)\,
     e^{-i\qp'\cdot\vc{\rho}_1}\,
     e^{-i(\qp-\qp')\cdot\vc{\rho}_2},
  \label{eq:v1_expand}\\
  \tilde\Phi_{x'y_3y_4}(\qp'-\qp,-\qp')
  &= \sum_{\vc{\rho}_3,\vc{\rho}_4}
     \Phi_{x'y_3y_4}(\vc{\rho}_3,\vc{\rho}_4)\,
     e^{-i(\qp'-\qp)\cdot\vc{\rho}_3}\,
     e^{+i\qp'\cdot\vc{\rho}_4}.
  \label{eq:v2_expand}
\end{align}
Collecting the four phase factors gives
\begin{equation}\label{eq:phases_combined}
  e^{-i\qp'\cdot\vc{\rho}_1}\,e^{-i(\qp-\qp')\cdot\vc{\rho}_2}\,
  e^{-i(\qp'-\qp)\cdot\vc{\rho}_3}\,e^{+i\qp'\cdot\vc{\rho}_4}
  =
  \underbrace{e^{+i\qp\cdot(\vc{\rho}_3-\vc{\rho}_2)}}_{\text{$\qp$-dep.}}
  \cdot
  \underbrace{e^{-i\qp'\cdot\vc{S}}}_{\text{$\qp'$-dep.}},
\end{equation}
with $\vc{S} = \vc{\rho}_1 - \vc{\rho}_2 + \vc{\rho}_3 - \vc{\rho}_4$.
The $\qp$-dependence is now a single phase
multiplying everything, and the $\qp'$-dependence is
in the second factor and in the two $g$-arguments. Pulling the
$\qp$-phase outside the $\qp'$-sum and applying the convolution
theorem to the $\qp'$-sum yields
\begin{equation}\label{eq:conv_result}
\begin{aligned}
  \frac{1}{N}\sum_{\qp'} e^{-i\qp'\cdot\vc{S}}\,
  G_{y_1y_4}(\omega';&\qp')\,G_{y_2y_3}(\omega-\omega';\qp-\qp')\\
  &= \sum_{\vR}
    g_{y_1y_4}(\omega';\vR-\vc{S})\,
    g_{y_2y_3}(\omega-\omega';\vR)\,e^{-i\qp\cdot\vR},
\end{aligned}
\end{equation}
where the shift by $\vc{S}$ in the first $g$ argument reflects the
multiplication by $e^{-i\qp'\cdot\vc{S}}$ in $\qp'$-space. A change
of variables $\vR' = \vR - (\vc{\rho}_3 - \vc{\rho}_2)$ combined
with the identity
\begin{equation}
  \vc{\rho}_3 - \vc{\rho}_2 - \vc{S}
  = \vc{\rho}_4 - \vc{\rho}_1
\end{equation}
cancels the $\qp$-dependent phase against the $e^{-i\qp\cdot\vR}$
factor produced by \cref{eq:conv_result}, leaving
\begin{equation}\label{eq:sigma_realspace}
  \Sigma_{xx'}^{\gtrless}(\omega;\qp)
  = \frac{i\hbar}{2}
    \sum_{\vR}\int\!d\tau\;
    \mathcal{B}_{xx'}(\tau;\vR)\,e^{+i\omega\tau}\,e^{-i\qp\cdot\vR},
\end{equation}
with the real-space bubble
\begin{equation}\label{eq:bubble_realspace}
\begin{aligned}
  \mathcal{B}_{xx'}(\tau;\vR)
  \equiv
  \!\!\sum_{\substack{\vc{\rho}_1,\ldots,\vc{\rho}_4\\y_1,\ldots,y_4}}\!\!
  &\Phi_{xy_1y_2}(\vc\rho_1,\vc\rho_2)\,
   \Phi_{x'y_3y_4}(\vc\rho_3,\vc\rho_4)\\
  &\times g_{y_1y_4}(\tau;\,\vR{+}\vc\rho_4{-}\vc\rho_1)\,
          g_{y_2y_3}(\tau;\,\vR{+}\vc\rho_3{-}\vc\rho_2).
\end{aligned}
\end{equation}
The structure of \cref{eq:bubble_realspace} replaces the explicit
$\qp'$-sum by a loop over $N_D \le N_\perp^{4}$ FC3-displacement
quadruples and an effective lattice $\mathcal{R}_\mathrm{eff}$ of
size $N_T^\mathrm{eff} = (4 N_\mathrm{cut}+1)^{2}$, the latter
following from the largest possible $\vR + \vc\rho_i - \vc\rho_j$
shift with $|\vc\rho_i|\le N_\mathrm{cut}$. The total cost is
\begin{equation}\label{eq:cost_real}
  \mathcal{W}_R = N_\tau\,N_D\,N_T^\mathrm{eff}\,\mathcal{C}_\mathrm{ic}.
\end{equation}
Comparing \cref{eq:cost_q,eq:cost_real}, the real-space approach
becomes favourable when $N_D N_T^\mathrm{eff} < N_q^{2}$. With
$N_\mathrm{cut} = 3$ and $N_\perp = 7$ (a typical FC3 support along
each transverse direction), $N_D \sim 7^{4} \approx 2400$ and
$N_T^\mathrm{eff} = 13^{2} \approx 170$, giving $N_D N_T^\mathrm{eff}
\sim 4\times 10^{5}$. For a $20\times 20$ transverse grid
($N_q = 400$, $N_q^{2} = 1.6\times 10^{5}$), the q-space form is
already cheaper, and the gap widens for finer Brillouin-zone
sampling.

% ----------------------------------------------------------------------
\subsection{Sparsity layers}
\label{ssec:phph_sparsity}
% ----------------------------------------------------------------------

Three block-level sparsity patterns reduce the dense ring cost
\cref{eq:cic_cost} by several orders of magnitude before any further
approximation is introduced. These mirror the sparsity exploitations
of the GW solver
\cite{vetsch_ab-initio_2025, deuschle_electron-electron_2025} and
follow from the same physical assumption -- a real-space cut-off on
inter-atomic interactions -- applied to the harmonic and cubic
parts of the lattice Hamiltonian. The cut-offs act at the granularity
of \emph{transport-cell blocks}: a block is either present in full
(stored as a dense $\NBS\times\NBS$ or $\NBS\times\NBS\times\NBS$
array) or absent. Within a present block the implementation uses
dense BLAS/einsum throughout
(\cref{ssec:phph_distribution,ssec:phph_dist_impl}); intra-block
atom-pair sparsity is \emph{not} exploited at compute time. CRS/COO
representations exist in the storage layer
(\texttt{DSDBSparse} in \cite{vetsch_ab-initio_2025}) only to (i)
record which block pairs are non-zero in the BT band and (ii) carry
the \texttt{stack}/\texttt{nnz} redistribution traffic; the inner
linear algebra always operates on the dense block view.

We denote by $r_{\mathrm{cut}}$ the physical real-space cut-off
radius applied to both IFC2 and IFC3 (in \AA), by $r_c$ the
corresponding bandwidth on the block-tridiagonal representation of
$G^{\gtrless}$ (the number of primitive blocks within $r_{\mathrm{cut}}$
along the transport direction), and by $N_\mathrm{c,x}$ the FC3
cut-off in transport-cell units along the transport direction
($N_\mathrm{c,x}\le 3$ typically). The transport-cell size $\NBS$
absorbs all the intra-cell degrees of freedom.

\paragraph{(a)~BT-band cut-off on $G^{\gtrless}$.} As in the
electronic GW problem, applying $r_{\mathrm{cut}}$ to the harmonic
dynamical matrix gives $G^{\gtrless}$ a block-banded structure of
width $r_c$ along the transport direction. Block entries $G_{KK'}$
are non-zero only for $|K-K'|\le r_c$, and each non-zero block is
stored densely as an $\NBS\times\NBS$ matrix. The total entry count
of $G$ per frequency is
\begin{equation}\label{eq:G_bb_nnz}
  \mathrm{nnz}(G)_\mathrm{BB}
  = \cO\!\bigl(\NB\,r_c\,\NBS^{2}\bigr)
  = \cO\!\bigl(N_{AO}\,r_c\,\NBS\bigr).
\end{equation}

\paragraph{(b)~Inter-block FC3 cut-off.} The decay of cubic force
constants with pair distance (\cref{sub:fc3_compression}) admits the
same real-space cut-off $r_{\mathrm{cut}}$ on $\Phi^{(3)}$. For
transport-cell labels $(I,K_1,K_2)$, $\tilde\Phi^{(3)}_{I,K_1,K_2}$
is non-zero (and stored as a dense $\NBS\times\NBS\times\NBS$
array) only when
\begin{equation}\label{eq:fc3_block_support}
  |I-K_1|, \quad |I-K_2|, \quad |K_1-K_2| \le N_\mathrm{c,x}.
\end{equation}
The block-pair index
\begin{equation}\label{eq:phph_pair_index}
  \mathcal{P}(I,J)
  = \bigl\{(K_1,K_2) :
      \Phi^{(3)}_{I,K_1,K_2}\;\text{and}\;
      \Phi^{(3)}_{J,K_2,K_1}\;\text{both present}\bigr\},
\end{equation}
encodes which inner block-pairs contribute to a given output
block-pair $(I,J)$. From \cref{eq:fc3_block_support} together with
the BT-band requirement on the output ($|I-J|\le 1$), the cardinality
of $\mathcal{P}(I,J)$ is bounded by
$|\mathcal{P}(I,J)| = \cO(N_\mathrm{c,x}^{2})$, with the proportionality
constant set by the lattice geometry. For a cubic lattice and
$N_\mathrm{c,x}=3$, $|\mathcal{P}(I,J)|\le 49$ in the worst case.
The partial Fourier transform $\Phi^{(3)} \to \tilde\Phi^{(3)}(\qp,\qp')$
over the transverse cells preserves block-sparsity along the transport direction unchanged.

\paragraph{(c)~Output BT band.} $\Sigma$ inherits the BT structure of
the upstream SCBA solver: only $(I,J)$ with $|I-J|\le 1$ are
retained, and each retained block is stored densely. The number of
distinct output block-pairs is $\cO(\NB)$, and the bubble decomposes
into a finite sum
\begin{equation}\label{eq:sigma_block_sum}
  \Sigma_{IJ}(\omega;\qp)
  = \frac{i\hbar}{2}
    \sum_{(K_1,K_2)\in\mathcal{P}(I,J)}
      \mathcal{B}_{IJ}^{K_1K_2}(\omega;\qp).
\end{equation}
This output truncation is consistent with the input truncation (a)
on $G$.

\paragraph{Per-pair kernel cost.} For each surviving block triple
$(I, K_1, K_2)$ and its conjugate $(J, K_2, K_1)$, the per-frequency
ring contraction
\cref{eq:phph_step_A,eq:phph_step_B,eq:phph_step_S} is evaluated as
three dense \texttt{einsum} contractions on dense
$\NBS^{3}$- and $\NBS^{2}$-shaped tensors. The intermediates
$A\in\mathbb{C}^{N_\tau\times\NBS\times\NBS\times\NBS}$ and
$B\in\mathbb{C}^{N_\tau\times\NBS\times\NBS\times\NBS}$ both have
$\cO(N_\tau\,\NBS^{3})$ entries; the dominant arithmetic is the
$B$-build at $\cO(N_\tau\,\NBS^{4})$ flops per pair (per $\tau$,
$\cO(\NBS^{4})$). The per-pair ring cost is therefore
\begin{equation}\label{eq:cic_block_dense}
  \mathcal{C}_\mathrm{ic}^\mathrm{block}
  = \cO\!\bigl(\NBS^{4}\bigr).
\end{equation}
This coincides with the dense-ring estimate \cref{eq:cic_cost} on a
single transport cell because there is no intra-block sparsity to
exploit; the savings relative to the naive global ring (which would
scan all of $\ndof = N_\mathrm{AO}$) come entirely from the
block-level cut-offs (a)--(c), not from intra-block sparsity. With
GPU-resident dense BLAS this kernel is compute-bound for
$\NBS\gtrsim 10^{2}$ and saturates the device's FP64 (or FP32)
throughput.

Combining \cref{eq:cic_block_dense} with the q-space cost
\cref{eq:cost_q} and the $\mathcal{P}$-cardinality bound
$|\mathcal{P}(I,J)|=\cO(N_\mathrm{c,x}^{2})$, the total q-space SSE
cost reduces to
\begin{equation}\label{eq:cost_phph_total}
  \mathcal{W}_\mathrm{phph}
  = \cO\!\bigl(N_\tau\,N_q^{2}\,\NB\,N_\mathrm{c,x}^{2}\,
               \NBS^{4}\bigr)
  = \cO\!\bigl(N_\tau\,N_q^{2}\,N_{AO}\,N_\mathrm{c,x}^{2}\,
               \NBS^{3}\bigr).
\end{equation}
The $\NBS^{3}$ scaling sets a hard upper bound on the
transport-cell size that can be handled at fixed compute budget and
motivates the FC3 decomposition of \cref{sec:phph_decomp}, which
replaces $\NBS^{3}$ by $R\,\NBS$ in the per-pair kernel.

\paragraph{Notation: intra-block dimension $C_a$ in later sections.}
The separable kernel of \cref{ssec:phph_separable} and the full
decomposition derivations of \cref{sec:phph_decomp} are written in
terms of a generic intra-block dimension $C_a$ -- the number of
atoms within $r_\mathrm{cut}$ of any reference atom, which would
control the per-block density if intra-block sparsity were
exploited by the inner kernel. The present quatrex implementation
runs dense BLAS within each block, so the effective value is
$C_a=\NBS$ and all cost comparisons in those sections collapse to
the block-dense baseline \cref{eq:cic_block_dense,eq:cost_phph_total}.
Retaining $C_a$ symbolically keeps the decomposition derivations
forward-compatible with a future implementation that exploits
intra-block sparsity (e.g.\ via block-CSR BLAS or atom-pair-sparse
einsum kernels); the operation counts in that hypothetical regime
would become $\NBS\to\NBS$, $\NBS^{2}\to\NBS\,C_a$,
$\NBS^{3}\to\NBS\,C_a^{2}$, $\NBS^{4}\to\NBS\,C_a^{3}$.

% ----------------------------------------------------------------------
\subsection{Separable FC3 approximation}
\label{ssec:phph_separable}
% ----------------------------------------------------------------------

The intra-block sparsity layer (d) reduces the per-pair ring cost
from $\NBS^{4}$ to $\NBS\,C_a^{3}$, but $C_a^{3}$ remains the
dominant arithmetic factor at production block sizes. The factor can
be reduced further by exploiting low-rank structure in the FC3
itself, replacing $\Phi$ by a sum-of-products approximation that
factorises the ring into a product of two independent bilinear
forms.

Treating $\Phi_{xy_1y_2}(\vc\rho_1,\vc\rho_2)$ as a matrix with row
index $(\vc\rho_1,y_1)$ and column index $(\vc\rho_2,y_2)$, a
rank-$R$ truncated SVD yields
\begin{equation}\label{eq:fc3_lowrank}
  \Phi_{xy_1y_2}(\vc\rho_1,\vc\rho_2)
  \approx \sum_{r=1}^{R}
    f^{x}_{y_1,r}(\vc\rho_1)\,
    h_{y_2,r}(\vc\rho_2),
\end{equation}
with $R \ll \NBS$, and analogously for the second FC3 factor with
factors $f',h'$ at rank $R$.  Substituting both into
\cref{eq:inner_contraction} gives the factorised ring
\begin{equation}\label{eq:sep_split}
  \mathcal{B}_{xx'}
  = \sum_{r,s}
    \underbrace{
      \Bigl[\sum_{y_1,y_4}f^{x}_{y_1,r}\,g_{y_1y_4}\,h'_{y_4,s}\Bigr]
    }_{P^{x}_{rs}}
    \cdot
    \underbrace{
      \Bigl[\sum_{y_2,y_3}h_{y_2,r}\,g_{y_2y_3}\,f'^{x'}_{y_3,s}\Bigr]
    }_{Q^{x'}_{rs}}.
\end{equation}
Naive evaluation of $P^{x}_{rs}$ as $R^{2}$ independent bilinears
costs $\cO(R^{2}\ndof^{2})$ per working point. Sharing the
matrix-vector products across $(r,s)$ pairs gives a two-stage
evaluation:
\begin{enumerate}
  \item For each $s$, compute the projected propagator
    $\vc{v}_s^{y_1} = \sum_{y_4}g_{y_1,y_4}\,h'_{y_4,s}$. With $g$
    sparse (nnz $\NBS\,C_a$ per block, banded with bandwidth $r_c$),
    each $\vc{v}_s$ costs $\cO(\ndof\,r_c\,C_a)$. Total over $s$:
    $\cO(R\,\ndof\,r_c\,C_a)$.
  \item For each $(r,s,x)$, compute the inner product
    $P^{x}_{rs} = \sum_{y_1}f^{x}_{y_1,r}\,\vc{v}_s^{y_1}$ at cost
    $\cO(C_a)$ per output (intra-block-sparse $f$). Total over
    $(r,s,x)$: $\cO(R^{2}\,\ndof\,C_a)$.
\end{enumerate}
The combined per-pair cost is
\begin{equation}\label{eq:cost_opt}
  \cO\!\bigl(R\,\ndof\,r_c\,C_a + R^{2}\,\ndof\,C_a\bigr),
\end{equation}
to be compared with $\cO(\ndof\,C_a^{3}/\NBS)$ from
\cref{eq:cic_block_dense} under the assumption that intra-block
sparsity is exploitable (cf.\ the forward note in
\cref{ssec:phph_sparsity}: in the current dense-kernel
implementation $C_a=\NBS$, and the comparison collapses to the
asymptotically equivalent $\cO(\NBS^{3})$ baseline). For
$R\ll C_a^{2}$ the separable form is the cheaper one.
Structurally richer decompositions (Tucker, CP, INDSCAL, symmetric
CP; \cref{sub:fc3_tucker,sub:fc3_cp,sub:fc3_indscal,sub:fc3_symcp})
plug into the same kernel with different factor matrices and the
same two-stage evaluation.

In the transverse-periodic setting, Fourier-transforming each
factor in \cref{eq:fc3_lowrank} turns the $\qp'$-sum at fixed
$(r,s)$ into a convolution evaluated by FFT. The dense
$\ndof^{2}\times N_q^{2}$ ring is replaced by $R^{2}$ pointwise
products of two FFT-convolved factors, each of which inherits the
two-stage assembly cost \cref{eq:cost_opt}.

% ----------------------------------------------------------------------
\subsection{Algebraic structure: FFT placement and evaluation order}
\label{ssec:phph_algebra}
% ----------------------------------------------------------------------

The bubble involves five elementary operations: forward FFT
$\mathcal{F}_\omega$ along $\omega$, inverse FFT
$\mathcal{F}^{-1}_\tau$ along $\tau$, contractions $\mathcal{C}_L$
and $\mathcal{C}_R$ with $\Phi_L$ and $\Phi_R$, and the pointwise
product $\mathcal{M}_\tau$ that closes the ring at fixed $\tau$.
Because $\mathcal{F}_\omega$ acts on the $\omega$-axis while
$\mathcal{C}_{L/R}$ acts only on index axes, the two commute,
\begin{equation}\label{eq:fft_contraction_commute}
  \mathcal{F}_\omega\circ\mathcal{C}_{L/R}
  = \mathcal{C}_{L/R}\circ\mathcal{F}_\omega,
\end{equation}
and we may choose the order of operations.

\paragraph{FFT placement.} Two orderings are possible.

\emph{(i) FFT-first.} Take the input $G^{\gtrless}(\omega)$ to the
$\tau$-domain once, and perform all FC3 contractions on $g(\tau)$.
The FFT input has $\mathrm{nnz}(G)$ entries (\cref{eq:G_bb_nnz}
times intra-block sparsity), each FFT'd independently along $\omega$
at cost $\cO(N_\omega\log N_\omega)$. Total FFT work:
\begin{equation}\label{eq:fft_cost_first}
  W_\mathrm{FFT}^\mathrm{first}
  = \cO\!\bigl(N_\omega\log N_\omega \cdot
               N_{AO}\,r_c\,C_a\bigr).
\end{equation}

\emph{(ii) Contract-first.} Perform all index contractions in the
$\omega$-domain to produce a four-index intermediate of FC3 shape,
then FFT the result. The intermediate has nnz proportional to
$\mathrm{nnz}(\Phi^{(3)})\sim N_{AO}\,N_\mathrm{c,x}^{2}\,C_a^{2}$,
each entry needs an FFT along $\omega$ at cost
$\cO(N_\omega\log N_\omega)$. Total FFT work:
\begin{equation}\label{eq:fft_cost_after}
  W_\mathrm{FFT}^\mathrm{after}
  = \cO\!\bigl(N_\omega\log N_\omega \cdot
               N_{AO}\,N_\mathrm{c,x}^{2}\,C_a^{2}\bigr).
\end{equation}

he FFT-first ordering is therefore the
clearly superior choice.

\paragraph{Per-pair kernel.} With the FFT-first ordering, the
balanced-tree (or, equivalently for one-rank evaluation, outside-in)
contraction at fixed $(I,J,K_1,K_2)\in\mathcal{P}$ takes the form
\begin{align}
  A_{w,a,c,d}
    &= \sum_{e}(\Phi_L)_{a,c,e}\,(g_b)_{w,e,d},
    \label{eq:phph_step_A}\\
  B_{w,a,b,d}
    &= \sum_{c}A_{w,a,c,d}\,(g_a)_{w,c,b},
    \label{eq:phph_step_B}\\
  \hat\Sigma_{w,a,J}
    &= \sum_{b,d}B_{w,a,b,d}\,(\Phi_R)_{J,d,b},
    \label{eq:phph_step_S}
\end{align}
where $w$ enumerates the padded $\tau$-grid of size
$N_\mathrm{fft}=2N_\omega-1$, $\Phi_L \equiv \Phi^{(3)}_{I,K_1,K_2}$
and $\Phi_R \equiv \Phi^{(3)}_{J,K_2,K_1}$ are the FC3 blocks for
the working triplet, and $g_a \equiv \mathcal{F}_\omega^{-1}\bigl[
G^{\gtrless}_{K_1\bullet}\bigr]$,
$g_b \equiv \mathcal{F}_\omega^{-1}\bigl[G^{\gtrless}_{K_2\bullet}\bigr]$
are the FFTs of the relevant blocks of $G$ along the transport
direction. The intermediates $A,B$ are stored only on the joint
sparse support of their inputs: $A$ on
$\mathrm{supp}(\Phi_L)\cap\mathrm{supp}(g_b)$ (which has nnz
$\cO(\NBS\,C_a^{2})$ per $\tau$ at fixed $(I,K_1,K_2)$), and
similarly for $B$.

An optimisation arises from sharing intermediates across FC3 triplets
that share an inner block. For fixed $K_1$, all
$(I,K_2)\in\mathcal{P}$ produce $A$-tensors that share the same
$g_a$-input. The $\mathcal{C}_L$ contraction can therefore be
batched into a single sparse einsum over the FC3 entries with
shared $K_1$.. Analogous sharing
applies across $K_2$.

% ----------------------------------------------------------------------
\subsection{Distribution and parallelisation}
\label{ssec:phph_distribution}
% ----------------------------------------------------------------------

The FFT-first pipeline factorises into three computational stages,
each of which is local in a different index. Let $P_E$ be the number of
ranks distributed along the $\omega$-axis, $P_S$ the number of ranks
distributed along the transport-direction blocks $I$ (the spatial
decomposition of \cite{vetsch_ab-initio_2025}), and $P=P_E P_S$ the
total rank count.

\paragraph{Stages.} We can split the computation in three stages,
\begin{align}
  g_{ij}(\tau)
    &= \mathcal{F}_{\omega\to\tau}\bigl[G_{ij}(\omega)\bigr]
    \qquad\text{(Stage~1: forward FFT, local in $(i,j)$)},
    \label{eq:phph_stage1}\\
  \sigma_{xx'}(\tau)
    &= \!\!\sum_{(K_1,K_2)\in\mathcal{P}(I,J)}\!\!
       \mathcal{B}_{IJ}^{K_1K_2}(\tau)
    \qquad\text{(Stage~2: ring contraction, local in $\tau$)},
    \label{eq:phph_stage2}\\
  \Sigma_{xx'}(\omega)
    &= \tfrac{i\hbar}{2}\,
       \mathcal{F}^{-1}_{\tau\to\omega}\bigl[\sigma_{xx'}(\tau)\bigr]
    \qquad\text{(Stage~3: inverse FFT, local in $(I,J)$)}.
    \label{eq:phph_stage3}
\end{align}
Stage~1 acts independently on each $(i,j)$ entry, Stage~2 acts independently at
each $\tau$ but couples block-diagonal $g$-blocks at different
$(K_1,K_2)$, and Stage~3 acts independently on each output $(I,J)$.

Following the GW solver, we use a block-banded sparse container with two
interconvertible distributions:
\begin{itemize}
  \item \texttt{stack}: each rank holds an $\omega$-slice (or
        $\tau$-slice after FFT) for the full BT band. The non-zero
        pattern of the BT band is replicated on every rank, the
        $\omega$-axis is partitioned across $P_E$.
  \item \texttt{nnz}: each rank holds the full $\omega$-axis for an
        $(i,j)$-slice of the BT band. The frequency axis is
        replicated, the non-zero pattern is partitioned across
        $P_E P_S = P$.
\end{itemize}
The \texttt{dtranspose} primitive converts between the two via a
balanced \texttt{Alltoall}. Stage~1 and Stage~3 are local in
\texttt{nnz}; Stage~2 is local in \texttt{stack}. Placing each stage in its preferred
distribution gives a four-redistribution pipeline:
\begin{enumerate}[label=(\roman*)]
  \item \texttt{stack}$\to$\texttt{nnz} on $G^{\gtrless}$.
  \item \texttt{nnz}: FFT $G\to g(\tau)$.
  \item \texttt{nnz}$\to$\texttt{stack} on $g(\tau)$.
  \item \texttt{stack}: per-$\tau$ ring contraction
        \cref{eq:phph_step_A,eq:phph_step_B,eq:phph_step_S} (Stage~2).
  \item \texttt{stack}$\to$\texttt{nnz} on $\sigma(\tau)$.
  \item \texttt{nnz}: IFFT $\sigma\to\Sigma$ (Stage~3).
  \item \texttt{nnz}$\to$\texttt{stack} on $\Sigma$ for hand-off
        to the next SCBA iteration.
\end{enumerate}
The four redistributions occur at steps (i), (iii), (v), (vii).

\paragraph{Per-rank cost derivation.} Each stage's cost follows
from the sparsity bookkeeping of \cref{ssec:phph_sparsity} together
with the distribution above.

\emph{Stage~1 (FFT, $G$ in \texttt{nnz})}. Memory: each rank holds
the full $\omega$-axis for its share of the BT-band non-zero blocks.
With $\mathrm{nnz}(G)$ per $\omega$ given by \cref{eq:G_bb_nnz}, the
per-rank footprint is
\begin{equation}
  M^{(1)}_\mathrm{rank}
  = \cO\!\Bigl(\frac{N_\omega \cdot N_{AO}\,r_c\,\NBS}{P}\Bigr).
\end{equation}
Computation: each non-zero $(I,J)$ block is FFT'd elementwise
along $\omega$; per-rank work is the same expression times
$\log N_\omega$:
\begin{equation}
  W^{(1)}_\mathrm{rank}
  = \cO\!\Bigl(\frac{N_\omega\log N_\omega \cdot N_{AO}\,r_c\,\NBS}{P}\Bigr).
\end{equation}
Communication: the \texttt{stack}$\to$\texttt{nnz} dtranspose has
the same volume as the memory footprint (this is the entire payload
that traverses the dtranspose).

\emph{Stage~2 (ring contraction, $g(\tau)$ in \texttt{stack})}. The
arithmetic at each working point is bounded by
\cref{eq:cic_block_dense}; summed over all working points and
distributed across all ranks,
\begin{equation}
  W^{(2)}_\mathrm{rank}
  = \cO\!\Bigl(\frac{N_q^{2}\,N_\tau\,N_{AO}\,N_\mathrm{c,x}^{2}\,\NBS^{3}}{P}\Bigr).
\end{equation}
Memory: dominated by the FC3 dictionary, which is partitioned along
its first leg $I$ across $P_S$ but not along $\omega$ (the FC3 is
$\omega$-independent). Each $P_S$ rank therefore holds
\begin{equation}
  M^{(2)}_\mathrm{rank}
  = \cO\!\Bigl(\frac{N_{AO}\,N_\mathrm{c,x}^{2}\,\NBS^{2}}{P_S}\Bigr).
\end{equation}
The $g(\tau)$ payload itself is in \texttt{stack} and adds an
$\cO(M^{(1)}_\mathrm{rank})$ contribution which is sub-leading at
production block sizes. Communication: the only Stage~2
communication is the spatial halo exchange discussed below.

\emph{Stage~3 (IFFT, $\sigma(\tau)$ in \texttt{nnz})}. The output
$\Sigma$ inherits the BT band of $G$: only $|I-J|\le 1$ blocks are
retained, each stored densely. The output entry count per $\omega$
is therefore $\cO(N_{AO}\,\NBS)$ (the BT-band factor $r_c$ of $G$
is replaced by the strict $|I-J|\le 1$ output band). Memory and
computation scale as
\begin{equation}
  M^{(3)}_\mathrm{rank}
    = \cO\!\Bigl(\frac{N_\omega\,N_{AO}\,\NBS}{P}\Bigr),
  \qquad
  W^{(3)}_\mathrm{rank}
    = \cO\!\Bigl(\frac{N_\omega\log N_\omega\,N_{AO}\,\NBS}{P}\Bigr).
\end{equation}

\paragraph{Spatial decomposition and halo.} When $P_S>1$ ranks
partition the transport-direction blocks, the FC3 cut-off implies
that for any output block $I\in\mathcal{I}_p$, only $g_{KK'}$ blocks
with $|K-I|, |K'-I|\le N_\mathrm{c,x}$ enter Stage~2. Each $P_S$
rank therefore needs only its assigned range $\mathcal{I}_p$ plus a
halo of width $N_\mathrm{c,x}$ blocks. The halo traffic per
$(P_E,P_S)$ rank is
\begin{equation}\label{eq:halo_volume}
  V_\mathrm{halo}^\mathrm{rank}
  = \cO\!\Bigl(\frac{N_\tau\,N_\mathrm{c,x}\,\NBS^{2}}{P_E}\Bigr),
\end{equation}
independent of $\NB$.

\paragraph{Aggregate per-rank cost.} The dominant cost contributions
balance across stages and ranks, giving
\begin{align}
  W^\mathrm{rank}_\mathrm{phph}
    &= \cO\!\Bigl(\frac{N_q^{2}\,N_\tau\,N_{AO}\,N_\mathrm{c,x}^{2}\,\NBS^{3}}{P}\Bigr),
    \label{eq:cost_phph_dist}\\
  V^\mathrm{rank}_\mathrm{phph}
    &= \cO\!\Bigl(\frac{N_\omega\,N_{AO}\,r_c\,\NBS}{P}\Bigr),
    \label{eq:vol_phph_dist}\\
  M^\mathrm{rank}_\mathrm{phph}
    &= \cO\!\Bigl(\frac{N_\omega\,N_{AO}\,r_c\,\NBS}{P}\Bigr)
       + \cO\!\Bigl(\frac{N_{AO}\,N_\mathrm{c,x}^{2}\,\NBS^{2}}{P_S}\Bigr).
    \label{eq:mem_phph_dist}
\end{align}
The two terms in \cref{eq:mem_phph_dist} reflect the two kinds of
data resident on a rank: the $\omega$-distributed $G$ and $\Sigma$
payloads (split over $P$) and the $\omega$-independent FC3
dictionary (split over $P_S$ only). Stage~2 dominates the
arithmetic; Stages~1 and 3 dominate the communication.

% ----------------------------------------------------------------------
\subsection{Memory and batching}
\label{ssec:phph_memory}
% ----------------------------------------------------------------------

The peak memory consumer in the per-pair kernel is the intermediate
$A$ or $B$ in \cref{eq:phph_step_A,eq:phph_step_B}, both of which
are dense $\NBS^{3}$-shaped tensors per $\tau$, i.e.\
$\cO(N_\mathrm{fft}\,\NBS^{3})$ entries across the padded grid.
At production block sizes ($\NBS\sim 2\times 10^{3}$,
$N_\omega\sim 10^{4}$), this is several TB per block triplet in
complex FP64 -- well beyond a single device. Two batching axes keep
the resident working set bounded:
\begin{itemize}
  \item The $\tau$-axis is point-wise within Stage~2 and can be
        batched without inducing additional communication. Batches
        of $N_b$ contiguous $\tau$-slices are sized adaptively from
        the free HBM at runtime, mirroring the GW polarisation
        kernel's batching strategy (\texttt{PCoulombScreening.compute}
        in \cite{vetsch_ab-initio_2025}).
  \item The output $(I,J)$-axis is batched analogously, since the
        contributions from different $(I,J)$ pairs are independent;
        the per-batch intermediate footprint is then
        $\cO(N_b\,\NBS^{3})$ per active $(I,J)$ pair.
\end{itemize}

The FC3 dictionary, with non-zero count
$\cO(\NB\,N_\mathrm{c,x}^{2}\,\NBS^{3})$ for block-dense storage,
is partitioned along its first leg $I$ across the $P_S$ spatial
ranks. With $\NBS\sim 2000$, $N_\mathrm{c,x}=3$ shells, the per-rank
FC3 footprint is hundreds of GB on a single device -- the
single-device regime is only practical for much smaller transport
cells ($\NBS\lesssim 300$). The compressed FC3 representations of
\cref{eq:fc3_lowrank,sub:fc3_compression} replace the per-block
$\NBS^{3}$ count by $R\,\NBS$ at moderate $R$ and integrate into
the same kernel and distribution scheme without further
modification; together with the GPU-resident dense BLAS that
already dominates Stage~2's arithmetic, they bring per-rank FC3
storage and per-pair arithmetic back to the few-GB / few-TFLOP
regime that fits a single HBM device.

% ----------------------------------------------------------------------
\subsection{Implemented distributed phonon-phonon SSE}
\label{ssec:phph_dist_impl}
% ----------------------------------------------------------------------

The current quatrex production implementation
(\texttt{quatrex.phonon.sse\_phonon\_phonon.SigmaPhononPhonon})
realises the dense-block kernel of \cref{eq:cic_block_dense} on
top of the distribution scheme of \cref{ssec:phph_distribution}.
It uses no FC3 decomposition (\cref{sec:phph_decomp} remains a
future optimisation), and it deliberately runs the contraction in
the \texttt{stack} distribution rather than the \texttt{nnz}
distribution used by the GW polarisation
\cite{vetsch_ab-initio_2025}: the phonon bubble couples block
triples through $\mathcal{P}(I,J)$ and so does not admit the same
per-nnz locality that makes the GW polarisation embarrassingly
parallel in \texttt{nnz}. The distribution exposes two levers:
\begin{itemize}
  \item ``\textbf{Frequency} $P_E$ via \texttt{comm.stack}.'' Each
        rank holds an $\omega$-slice of every BT-band $G$ and
        $\Sigma$ block. Inside the bubble, the diagonal blocks
        $g_{KK}(\omega)$ are all-gathered along $\omega$ across
        \texttt{comm.stack} into a full-axis tensor of shape
        $(N_\omega, \NBS, \NBS)$. Only diagonal blocks enter the
        contraction (\cref{eq:phph_pair_index} restricts to
        $G_{K_1K_1}\,G_{K_2K_2}$), so off-diagonal BT-band $G$
        blocks are never communicated.
  \item ``\textbf{Spatial} $P_S$ via \texttt{comm.block}.'' Each
        rank owns a contiguous global-block range
        $[I_\mathrm{lo}, I_\mathrm{hi})$ determined by
        \texttt{block\_section\_offsets}. The outer
        $(I,J)\in\mathcal{P}$ loop is filtered to $I\in
        [I_\mathrm{lo}, I_\mathrm{hi})$ so each rank contributes
        only the output blocks it owns; cross-rank diagonal blocks
        are made available through an additional
        \texttt{comm.block.all\_gather} of the locally-owned
        diagonals (requires uniform block sizes, which
        \texttt{SCBAData} enforces by default).
\end{itemize}

\paragraph{Pipeline.} A single SCBA iteration invokes the
implemented routine in three phases, all stack-resident:
\begin{enumerate}[label=(\roman*)]
  \item \emph{State restoration contract.} On entry the routine
        records each buffer's incoming distribution state and, if
        necessary, dtranspose's everything to \texttt{stack}. On
        exit (\texttt{finally} block) the incoming state is
        restored on every buffer. This makes the routine composable
        with the SCBA loop's \texttt{stack}$\to$\texttt{nnz}
        dtranspose: the GW polarisation runs in \texttt{nnz} and
        the phph SSE returns control with everything still in
        \texttt{nnz}, exactly as the surrounding driver expects.
  \item \emph{Diagonal-only gather.} For each local block $K$, the
        $\omega$-axis of $g_{KK}^{<,>}$ is all-gathered across
        \texttt{comm.stack} into a full-axis array. When
        $P_S>1$ a subsequent \texttt{comm.block.all\_gather} of the
        stacked local diagonals produces, on every rank, the full
        global $\{g_{KK}\}_{K=0}^{N_B-1}$ dictionary at full
        $\omega$.
  \item \emph{Per-rank bubble + write-back.} For each
        $(I,J)\in\mathcal{P}$ with $I\in[I_\mathrm{lo},
        I_\mathrm{hi})$ and $|I-J|\le 1$, the rank walks the
        precomputed block-pair index
        $\mathcal{P}(I,J)$, accumulates
        $\Sigma^{<,>}_{IJ}(\omega)$ via the dense
        \texttt{bubble\_dense} kernel (three \texttt{einsum}s on
        dense $\NBS$-shaped tensors plus an FFT along $\omega$),
        forms $\Sigma^R$ from the bosonic Hilbert transform
        $\Sigma^R = \tfrac{1}{2}(\Sigma^>{-}\Sigma^<)
        + \tfrac{i}{2}\mathcal{H}\{\Sigma^>{-}\Sigma^<\}$, and
        writes the $[e_\mathrm{lo}{:}e_\mathrm{hi}]$
        $\omega$-slice -- where
        $e_\mathrm{lo}=\sum_{r<\mathtt{stack.rank}}
        \mathtt{stack\_section\_sizes}_r$ -- into the
        local-block-indexed output buffer
        $\Sigma_{\bullet}.\mathtt{blocks}[I_\mathrm{loc},
        J_\mathrm{loc}]$.
\end{enumerate}

\paragraph{Per-rank cost.} The dominant arithmetic is the dense
ring of \cref{eq:cic_block_dense} summed over the rank's owned
outputs, giving
\begin{equation}\label{eq:cost_phph_impl}
  W^\mathrm{rank}_{\mathrm{phph,impl}}
  = \cO\!\Bigl(
      \frac{N_\omega\,\log N_\omega \cdot \NB\,N_\mathrm{c,x}^{2}\,\NBS^{4}}{P_S}
    \Bigr),
\end{equation}
where the $P_E$ factor does \emph{not} appear because each
\texttt{comm.stack} rank still computes the bubble at full
$\omega$ on its assigned $(I,J)$ outputs (this is the cost of not
distributing along $\omega$; recovering it would require a
distributed 1-D FFT, deferred). Communication is dominated by the
diagonal-block gather:
\begin{equation}\label{eq:vol_phph_impl}
  V^\mathrm{rank}_{\mathrm{phph,impl}}
  = \cO\!\Bigl(\frac{N_\omega\,\NB\,\NBS^{2}}{P}\Bigr),
\end{equation}
i.e.\ the per-$\omega$ payload of $N_B$ diagonal blocks of size
$\NBS\times\NBS$, divided across the $P=P_EP_S$ ranks. The FC3
dictionary is partitioned along $I$ across $P_S$ only:
\begin{equation}\label{eq:mem_phph_impl}
  M^\mathrm{rank}_{\mathrm{phph,impl}}
  = \cO\!\Bigl(\frac{N_\omega\,\NB\,\NBS^{2}}{P}\Bigr)
  + \cO\!\Bigl(\frac{\NB\,N_\mathrm{c,x}^{2}\,\NBS^{3}}{P_S}\Bigr).
\end{equation}

\paragraph{Validation.} The implementation is pinned against the
single-block dense reference of
\texttt{phonon.solver.dense.compute\_phph\_self\_energy\_finite}
(see \texttt{tests/quatrex/phonon/test\_sse\_phonon\_phonon.py})
in two regimes: (i) one transport cell, where the production
\texttt{\_bubble\_block} reproduces the dense reference to
$<10^{-10}$ relative Frobenius; (ii) multi-block ($N_B\ge 3$),
where an independently written reference assembles the BTD
$\Sigma_{IJ}$ output from the same FC3 dictionary and the same
diagonal $G$ blocks, again matching to $<10^{-10}$. A second test
pins the state-restoration contract.

\paragraph{Composition with the Interaction registry.}
The implemented bubble is registered as
\texttt{PhononPhononInteraction}, one of several concrete
\texttt{Interaction} instances iterated by the SCBA driver
(\texttt{quatrex.core.scba.SCBA.\_compute\_interactions}). The
registry is the architectural seam for adding a coupled
electron-phonon SSE: a future
\texttt{ElectronPhononInteraction} reads from both the electron
and phonon subsystems' Green's functions and writes its
contribution into the appropriate $\sigma_*$ buffer, using the
same \texttt{comm.stack}/\texttt{comm.block} grid and the same
\texttt{dtranspose} machinery as the existing interactions, with
no further driver-side changes required. In particular, the
state-restoration contract above guarantees that
\texttt{PhononPhononInteraction.compute} can be sequenced before
or after a future \texttt{ElectronPhononInteraction.compute}
without spurious distribution-state transitions between them.

% ----------------------------------------------------------------------
\section{Phonon-Phonon SSE with decomposed FC3}
\label{sec:phph_decomp}
% ----------------------------------------------------------------------

The sparsity layers of \cref{ssec:phph_sparsity} reduce the per-pair
ring contraction from $\cO(\NBS^{4})$ to $\cO(\NBS\,C_a^{3})$. The
factor $C_a^{3}$ remains the dominant const in
\cref{eq:cost_phph_total}. We now show that any decomposition of the FC3
in \cref{sub:fc3_compression} reduces this further to
$\cO(R\,\NBS\,r_c\,C_a + R^{2}\NBS)$ by collapsing the four-cycle of
the ring contraction into a single Gram matrix on the propagator.
The reduction is independent of which decomposition is used in the
sense that all CP-family decompositions (CP, INDSCAL, Waring, PCP) produce
the same kernel structure. The symmetry constraints distinguishing
them affect only which factor vectors enter the Gram product and how
much it can be symmetrised.

We derive this under two assumptions that hold for the SSE of a single
SCBA iteration on a single FC3 dataset. First, both vertices of the
sunset diagram (\cref{fig:sunset_diagram}) are evaluated against the
same $\Phi^{(3)}$: $\Phi_L \equiv \Phi_R \equiv \Phi$. Second, the
two internal propagator lines are restrictions of the same Green's
function $g(\tau)$ to (possibly different) block-row supports.

The section is organised as follows.
\Cref{ssec:phph_decomp_setup} fixes the notation and the abstract
separable form. \Cref{ssec:phph_decomp_kernel} derives the factorised
kernel and identifies the algebraic simplification that distinguishes
the symmetric ans\"atze from plain CP.
\Cref{ssec:phph_decomp_ansatze} specialises to each ansatz of
\cref{sub:fc3_compression}. \Cref{ssec:phph_decomp_blocks} traces
through the three block-sparsity layers of \cref{ssec:phph_sparsity}
under decomposition. \Cref{ssec:phph_decomp_compute} gives a concrete
three-step algorithm and its cost. \Cref{ssec:phph_decomp_qspace}
addresses the partial Fourier transform.
\Cref{ssec:phph_decomp_dist} reworks the distributed pipeline of
\cref{ssec:phph_distribution}.
\Cref{ssec:phph_decomp_asr,ssec:phph_decomp_error} cover ASR and
observable-error propagation. \Cref{ssec:phph_decomp_summary}
collects the asymptotic results.

% ----------------------------------------------------------------------
\subsection{Setup}
\label{ssec:phph_decomp_setup}
% ----------------------------------------------------------------------

Every decomposition considered in \cref{sub:fc3_compression} writes the FC3 vertex as a sum of
separable terms. The most general form is
\begin{equation}\label{eq:fc3_separable_general}
  \Phi_{x y_1 y_2}
  \;\approx\;
  \sum_{r=1}^{R} \lambda_r\;\alpha_r(x)\;\beta_r(y_1)\;\gamma_r(y_2),
\end{equation}
with weights $\lambda_r\in\mathbb{R}$ and factor vectors
$\alpha_r,\beta_r,\gamma_r\in\mathbb{R}^{\ndof}$. The symmetry constraints are
\begin{equation}\label{eq:symmetry_constraints}
\begin{aligned}
  \text{CP:}      \quad& \text{no constraint,}\\
  \text{INDSCAL:} \quad& \beta_r=\gamma_r,\\
  \text{Waring:}  \quad& \alpha_r=\beta_r=\gamma_r,\\
  \text{Tucker:}  \quad& \text{$\beta,\gamma$ taken from a common $R_2$-dim basis,}\\
                   & \text{$\alpha$ from a separate $R_1$-dim basis, plus a small core.}
\end{aligned}
\end{equation}
These constraints reduce parameter count
(\cref{eq:fc3_matsvd_params,eq:fc3_tucker_params,eq:fc3_cp_params,eq:fc3_indscal_params,eq:fc3_waring_params,eq:fc3_pcp_params})
and enforce permutation symmetry of
the reconstructed tensor in $(y_1,y_2)$ (INDSCAL, Waring) or all
three legs (Waring, PCP). PCP is treated separately in
\cref{ssec:phph_decomp_ansatze} because its $6$-term explicit
symmetrisation interacts non-trivially with the kernel.

We collect the factor vectors as columns of three matrices
$A,B,C \in \mathbb{R}^{\ndof\times R}$ with
$A_{x,r}\!=\!\alpha_r(x)$, $B_{y,r}\!=\!\beta_r(y)$,
$C_{y,r}\!=\!\gamma_r(y)$, and the weights as a diagonal matrix
$\Lambda = \diag(\lambda_1,\dots,\lambda_R)$. Under the constraints
\cref{eq:symmetry_constraints}, $B=C$ for INDSCAL and $A=B=C\equiv U$
for Waring. The propagator $g\equiv g(\tau)\in\mathbb{R}^{\ndof\times\ndof}$
is the time-domain Green's function obtained from Stage~1 of
\cref{ssec:phph_distribution}. We will not yet distinguish lesser,
greater, or retarded components, the derivation applies unchanged to
all three.

% ----------------------------------------------------------------------
\subsection{Factorised kernel}
\label{ssec:phph_decomp_kernel}
% ----------------------------------------------------------------------

The per-pair ring \cref{eq:inner_contraction} is
\begin{equation}
  \mathcal{B}_{xx'}
  = \sum_{y_1 y_2 y_3 y_4}
    \Phi(x,y_1,y_2)\,g(y_1,y_4)\,g(y_2,y_3)\,
    \Phi(x',y_3,y_4).
\end{equation}
Substituting \cref{eq:fc3_separable_general} for both copies of
$\Phi$, factoring the $g$-independent parts, and collecting yields
\begin{align}
  \label{eq:bubble_factorised}
  \mathcal{B}_{xx'}
  &= \sum_{r,s=1}^{R} \lambda_r\,\lambda_s\;
      \alpha_r(x)\,\alpha_s(x')\;
      P_{rs}\,Q_{rs},
  \\
  \label{eq:gram_P}
  P_{rs} &\equiv \beta_r^{\!\top} g\,\gamma_s \;=\; (B^{\!\top} g C)_{rs},
  \\
  \label{eq:gram_Q}
  Q_{rs} &\equiv \gamma_r^{\!\top} g\,\beta_s \;=\; (C^{\!\top} g B)_{rs}.
\end{align}
The four-cycle in the original ring collapses to two independent
two-cycles. Each two-cycle is a scalar bilinear form on the
propagator. The full $\NBS\times\NBS$ propagator entered
\cref{eq:inner_contraction} four times through tensor contractions. In
\cref{eq:bubble_factorised,eq:gram_P,eq:gram_Q}, it enters twice as the
kernel of a quadratic form.

The matrices $P$ and $Q$ defined in \cref{eq:gram_P,eq:gram_Q} are
not independent. They satisfy the algebraic identity
\begin{equation}\label{eq:Q_transpose}
  Q \;=\; C^{\!\top} g\,B \;=\; (B^{\!\top} g^{\!\top} C)^{\!\top}.
\end{equation}

\emph{INDSCAL ($B=C$).} The two bilinear forms are
\begin{equation}\label{eq:P_indscal}
  P_{rs} \;=\; \beta_r^{\!\top} g\,\beta_s
         \;=\; (B^{\!\top} g B)_{rs}
         \;=\; Q_{rs},
\end{equation}
identical entry-by-entry. The result holds without any symmetry
assumption on $g$.

\emph{Waring ($A=B=C=U$).} Same as INDSCAL, with the additional
property that the assembly factor $\alpha$ also reduces to $U$:
\begin{equation}\label{eq:bubble_waring}
  \mathcal{B}_{xx'}^{\mathrm{Waring}}
  = \sum_{r,s} \lambda_r \lambda_s\,U_{x,r}\,U_{x',s}\,
      \bigl[(U^{\!\top} g\,U)_{rs}\bigr]^{2}.
\end{equation}
The bilinear form is squared elementwise; the entire kernel is
parametrised by a single $\ndof\times R$ matrix $U$ and the weights
$\lambda$.

For CP without symmetry, $P\ne Q$ in general.

\paragraph{Unified kernel.}
Under the simplifications above, the bubble \cref{eq:bubble_factorised}
can be written as a single GEMM--Hadamard--GEMM sequence:
\begin{align}
  \label{eq:kernel_step1}
  V       &= g\,W,                          \\
  \label{eq:kernel_step2}
  P       &= W^{\!\top} V,                     \\
  \label{eq:kernel_step3}
  H_{rs}  &= \lambda_r\,\lambda_s\,P_{rs}\,P^{\star}_{rs}, \\
  \label{eq:kernel_step4}
  \Sigma  &= A\,H\,A^{\!\top},
\end{align}
where the choice of $W$ and $P^{\star}$ depends on the ansatz:
\begin{equation}\label{eq:kernel_choices}
\begin{aligned}
  \text{Waring:}\;\; & W = U,\;\; P^\star = P, \;\; A = U, \\
  \text{INDSCAL:}\;\; & W = B,\;\; P^\star = P, \;\; A = A_{\mathrm{INDSCAL}}, \\
  \text{CP:}\;\; & W = (B,C),\;\; P^\star = Q, \;\; A = A_{\mathrm{CP}}.
\end{aligned}
\end{equation}
For CP, $W=(B,C)$ denotes that two SpMMs are needed (one against $B$,
one against $C$) to obtain $V_B=gB$ and $V_C=gC$.INDSCAL and Waring need only
one SpMM. This is the dominant cost saving from imposing internal
symmetry on the decomposition.

% ----------------------------------------------------------------------
\subsection{Specialisation to each ansatz}
\label{ssec:phph_decomp_ansatze}
% ----------------------------------------------------------------------

The CP-family ans\"atze all fit the unified kernel
\cref{eq:kernel_step1,eq:kernel_step2,eq:kernel_step3,eq:kernel_step4}
with the choices \cref{eq:kernel_choices}. Tucker and the
matricisation SVD require minor modifications, given below.

\paragraph{CP.}
$\Phi_{xy_1y_2}^{\mathrm{CP}} = \sum_{r} \lambda_r\,a_r(x)\,b_r(y_1)\,c_r(y_2)$
with $A,B,C$ independent. The kernel requires both $V_B = gB$ and
$V_C = gC$; $P = B^{\!\top}V_C$ and (assuming $g=g^{\!\top}$) $Q=P^{\!\top}$,
giving $H_{rs}=\lambda_r\lambda_s P_{rs}P_{sr}$. The reconstructed
tensor is generically not $S_2$-symmetric in $(y_1,y_2)$; this
asymmetry survives into $\Sigma$ unless restored in post-processing.

\paragraph{INDSCAL.}
$\Phi_{xy_1y_2}^{\mathrm{INDSCAL}} = \sum_{r}\,d_r(x)\,v_r(y_1)\,v_r(y_2)$
with $B=C\equiv V$ and $A_\mathrm{INDSCAL} = D = [d_1\,\dots\,d_R]$.
The kernel computes $V_B = gV$, $P=V^{\!\top}V_B$ (a single GEMM-like
product), and assembles via $\Sigma = D\,(P\circ P)\,D^{\!\top}$ with
the elementwise square $(P\circ P)_{rs}=\lambda_r\lambda_s P_{rs}^{2}$.
The reconstructed tensor is exactly $S_2$-symmetric, so the SSE
inherits exact $\Sigma_{xx'}=\Sigma_{x'x}$ at the algorithmic level.

\paragraph{Waring.}
$\Phi_{xy_1y_2}^{\mathrm{Waring}} = \sum_{r}\lambda_r\,u_r(p(x))\,u_r(y_1)\,u_r(y_2)$,
where $p$ is the supercell lift of \cref{sub:fc3_symcp}. With
$A=B=C=U$, the kernel reduces further: a single matrix $U$
parametrises both the SpMM input and the output assembly. The bubble
is symmetric in $(r,s)$ by construction, so $P$ is symmetric
($P=P^{\!\top}$) and may be computed as a SYRK (upper-triangular GEMM)
at half the cost of the general INDSCAL GEMM.

\paragraph{PCP.}
$\Phi^{\mathrm{PCP}}_{xy_1y_2} = \sum_{\xi}\tfrac{\lambda_\xi}{6}
\sum_{\sigma\in S_3}
u^{\xi}_{\sigma_1}(p(x))\,u^{\xi}_{\sigma_2}(y_1)\,u^{\xi}_{\sigma_3}(y_2)$.
Each $\sigma\in S_3$ contributes a CP-shaped term, so PCP is
algebraically equivalent to CP at rank $6R$ with a specific structure
on the $6R$ triplets. The kernel
\cref{eq:kernel_step1,eq:kernel_step2,eq:kernel_step3,eq:kernel_step4}
applies with $R \to 6R$. For SSE evaluation, PCP is dominated by
Waring at matched fitting cost; we include it only for completeness.

\paragraph{Tucker.}
$\Phi_{xy_1y_2}^{\mathrm{Tucker}}=\sum_{pqr}G_{pqr}A_{xp}B_{y_1 q}B_{y_2 r}$
with $A\in\mathbb{R}^{\ndof\times R_1}$, $B\in\mathbb{R}^{\ndof\times R_2}$
columnwise orthonormal and a core $G\in\mathbb{R}^{R_1\times R_2\times R_2}$
with $G_{pqr}=G_{prq}$. The factorised kernel becomes matrix-valued:
\begin{align}
  \label{eq:tucker_kernel_step1}
  V_B    &= g\,B \in\mathbb{R}^{\ndof\times R_2},\\
  \label{eq:tucker_kernel_step2}
  M      &= B^{\!\top} V_B \in\mathbb{R}^{R_2\times R_2},\\
  \label{eq:tucker_kernel_step3}
  H_{pp'} &= \sum_{q r q' r'} G_{p q r}\,G_{p' q' r'}\,M_{q r'}\,M_{r q'},\\
  \label{eq:tucker_kernel_step4}
  \Sigma &= A\,H\,A^{\!\top},
\end{align}
The bilinear form is a single $R_2\times R_2$ Gram matrix $M$ (a
single SpMM-GEMM); the core contraction
\cref{eq:tucker_kernel_step3} is purely propagator-independent and
has cost $\cO(R_1^{2}R_2^{3})$ after $S_2$ symmetrisation. The
propagator-side cost is $\cO(R_2\,\NBS\,r_c\,C_a + R_2^{2}\NBS)$,
identical in form to the CP family with $R\to R_2$.

% ----------------------------------------------------------------------
\subsection{Block sparsity}
\label{ssec:phph_decomp_blocks}
% ----------------------------------------------------------------------

The unified kernel
\cref{eq:kernel_step1,eq:kernel_step2,eq:kernel_step3,eq:kernel_step4}
contracts $g$ against a dense factor matrix $W$ in step
\cref{eq:kernel_step1} and contracts $H$ against $A$ in step
\cref{eq:kernel_step4}. Two of the four block-sparsity layers of
\cref{ssec:phph_sparsity} survive the decomposition unchanged; one is
absorbed into the factor vectors; one becomes the support pattern of
$P$.

\paragraph{(a) BT-band of $g$ — preserved.}
The cut-off $r_c$ on the dynamical matrix gives $g(\tau)$
a block-banded structure of bandwidth $r_c$ along the transport
direction (\cref{eq:G_bb_nnz}). The SpMM $V=g\,W$ is therefore a
block-banded matrix times a dense $\ndof\times R$ matrix. The
non-zero pattern of $V$ is dictated by $g$: for any block $I$ on the
transport axis, $V_I = \sum_{|J-I|\le r_c}\,g_{IJ}\,W_J$, a band of
width $2r_c+1$ in the input. The non-zero count of $V$ is
\begin{equation}\label{eq:V_nnz}
  \mathrm{nnz}(V) \;=\; R\cdot\mathrm{nnz}_{\text{cols}}(g)
                 \;=\; \cO\!\bigl(R\,\NBS\,\NB\bigr)
                 \;=\; \cO\!\bigl(R\,\ndof\bigr),
\end{equation}
provided $W$ is dense (factor vectors with full support along
transport). The cost of the SpMM is
$\cO(R\cdot\mathrm{nnz}(g))=\cO(R\,\ndof\,r_c\,C_a)$.

\paragraph{(b) Intra-block atom-pair sparsity of $g$ — preserved.}
Within each $\NBS\times\NBS$ block, $g_{IJ}$ has nominal density
$C_a/\NBS$ (with $C_a=\NBS$ in the present dense-block kernel; see
the forward note in \cref{ssec:phph_sparsity}). The block-level SpMM
$V_I=\sum_J g_{IJ}\,W_J$ is therefore a sparse--dense block product
with cost
$\cO(R\,\NBS\,C_a)$ per source block $J$, summed over the $\cO(r_c)$
blocks within bandwidth, giving $\cO(R\,\NBS\,r_c\,C_a)$ per output
block and $\cO(R\,\ndof\,r_c\,C_a)$ for the full SpMM.

\paragraph{(c) FC3 inter-block support — absorbed.}
The FC3 cut-off
\cref{eq:fc3_block_support} restricts $\Phi^{(3)}_{I,K_1,K_2}$ to
block triples with mutual transport-direction distance
$\le N_\mathrm{c,x}$. After decomposition, the FC3 cut-off is
absorbed into the support of the factor vectors. Specifically, for
a CP fit of the full FC3 tensor, each rank-one component
$\lambda_r\,a_r\otimes b_r\otimes c_r$ has full support
$\mathrm{supp}(a_r)\times\mathrm{supp}(b_r)\times\mathrm{supp}(c_r)$;
all such supports must lie within the FC3 cut-off, so each individual
factor vector has support within $\cO(N_\mathrm{c,x}+\mathrm{spread})$
transport blocks. The total non-zero count of $W$ is therefore
\begin{equation}\label{eq:W_nnz}
  \mathrm{nnz}(W) \;=\; \cO\!\bigl(R\,\NBS\,N_\mathrm{c,x}\bigr)
                 \;=\; \cO\!\bigl(R\,\ndof\,N_\mathrm{c,x}/\NB\bigr),
\end{equation}
substantially smaller than the dense $R\,\ndof$ when
$N_\mathrm{c,x}\ll\NB$ (true for any transport device longer than
the FC3 cut-off).

In practice, an unconstrained CP-ALS fit produces factor vectors
without compact support: the rotational gauge of CP allows globally
delocalised factors to represent locally supported tensors. To
preserve the FC3 cut-off as factor-vector support, one of three
strategies is used: (i) initialise CP-ALS from a spatially localised
seed (e.g.\ the slice-wise SVD initialisation of
\cref{sub:fc3_indscal}); (ii) regularise with an $\ell_1$ or
group-sparse penalty as in
\cite{papalexakis_parcube_2012, kim_sparse_2007}; or (iii) work in
the CANDELINC framework of \cref{sub:fc3_other} with a basis that
respects the FC3 cut-off. All three preserve \cref{eq:W_nnz} at
matched accuracy; the choice is empirical and beyond the scope of
the SSE kernel.

\paragraph{(d) Output BT band — preserved by output projection.}
The output $\Sigma$ lives on the same BT band as $G$, with
$|I-J|\le 1$. The assembly step \cref{eq:kernel_step4} is therefore
restricted to this band. \emph{Crucially}, this restriction is
applied at the output, not at the intermediate $H$: the $R\times R$
matrix $H$ has dense support and does not inherit the band structure
of $\Sigma$. We compute $H$ once per $\tau$ and project to the BT
band when constructing $\Sigma$.

\paragraph{Effective sparsity of $P$.}
The Gram matrix $P$ has entries $P_{rs} = \beta_r^{\!\top} g\,\gamma_s$.
Two effects can render $P$ sparse:
\begin{itemize}
  \item \emph{Disjoint factor supports.} If
    $\mathrm{supp}(\beta_r)\cap N_{r_c}(\mathrm{supp}(\gamma_s))=\emptyset$
    (where $N_{r_c}$ is the $r_c$-block neighbourhood induced by $g$'s
    bandwidth), then $P_{rs}=0$ exactly.
  \item \emph{Approximate orthogonality.} For Waring with
    fitted $u_r$ that are approximately orthonormal in the metric
    $g$, the off-diagonal entries of $P$ are small. This is a
    refit-dependent property and we do not assume it.
\end{itemize}
The first effect can be substantial: for FC3 with $N_\mathrm{c,x}=3$
and a device of $\NB$ transport blocks, the fraction of $P_{rs}$
pairs with overlapping supports scales as
$N_\mathrm{c,x}/\NB$ — between $10^{-2}$ and $10^{-1}$ at production
device sizes. Exploiting this sparsity in step
\cref{eq:kernel_step2} converts the dense $R^{2}\NBS$ GEMM into a
sparse--sparse product at cost $\cO(R^{2}N_\mathrm{c,x})$, dominated
by the SpMM cost.

% ----------------------------------------------------------------------
\subsection{Computation strategy}
\label{ssec:phph_decomp_compute}
% ----------------------------------------------------------------------

The per-$\tau$ kernel
\cref{eq:kernel_step1,eq:kernel_step2,eq:kernel_step3,eq:kernel_step4}
maps onto three numerical primitives, all standard in modern linear
algebra libraries:
\begin{enumerate}
  \item \emph{Step \cref{eq:kernel_step1}: batched sparse SpMM.}
        $V = g\,W$, where $g$ is block-banded sparse and $W$ is
        $\ndof\times R$ with FC3-induced support \cref{eq:W_nnz}.
        Implemented as a sparse matrix times dense matrix product
        with $R$ right-hand sides batched into a single SpMM call.
        Cost: $\cO(R\,\ndof\,r_c\,C_a)$ flops, memory traffic
        $\cO(\mathrm{nnz}(g) + R\,\ndof)$. With $R$ batched RHSs,
        the arithmetic intensity is $\cO(R)$ flops/byte, placing the
        kernel firmly in the compute-bound regime of modern
        accelerators for $R\gtrsim 32$.
  \item \emph{Step \cref{eq:kernel_step2}: small dense GEMM.}
        $P = W^{\!\top}V$, a dense $R\times\ndof$ times $\ndof\times R$
        product yielding an $R\times R$ matrix.
        For Waring, $P$ is symmetric and the product reduces to
        SYRK. Cost: $\cO(R^{2}\,\ndof)$ flops. With $R^{2}$ output
        entries and $\cO(R\,\ndof)$ input non-zeros, the arithmetic
        intensity is $\cO(\ndof/R)$ flops/byte, comfortably
        compute-bound.
  \item \emph{Step \cref{eq:kernel_step4}: sparse-output sandwich.}
        $\Sigma = A\,H\,A^{\!\top}$ restricted to the output BT band
        and any available intra-block atom-pair sparsity
        (with $C_a=\NBS$ in the present dense-block kernel; see the
        forward note in \cref{ssec:phph_sparsity}).
        Implemented as: (i) precompute $T=AH$, a dense $\ndof\times R$
        matrix, cost $\cO(R^{2}\ndof)$; (ii) for each non-zero
        entry $(x,x')$ in the output support, compute
        $\Sigma_{xx'}=T_{x,:}\cdot A_{x',:}^{\!\top}$, cost
        $\cO(R)$. Total: $\cO(R^{2}\ndof + R\,\ndof\,C_a)$.
\end{enumerate}
The Hadamard step \cref{eq:kernel_step3} has cost $\cO(R^{2})$ and
is negligible.

\paragraph{Per-$\tau$ cost.}
Summing the dominant terms,
\begin{equation}\label{eq:per_tau_cost}
  \mathcal{C}_{\tau}^{\mathrm{decomp}}
  \;=\;
  \underbrace{\cO\!\bigl(R\,\ndof\,r_c\,C_a\bigr)}_{\text{SpMM}}
  \;+\;
  \underbrace{\cO\!\bigl(R^{2}\,\ndof\bigr)}_{\text{Gram + assembly}}.
\end{equation}
This is to be compared with the sparse-ring per-$\tau$ cost
\cref{eq:cost_phph_total},
\begin{equation}\label{eq:per_tau_cost_sparse}
  \mathcal{C}_{\tau}^{\mathrm{sparse}}
  \;=\;
  \cO\!\bigl(\ndof\,N_\mathrm{c,x}^{2}\,C_a^{3}\bigr).
\end{equation}
At Si-nanoribbon parameters ($C_a\sim 30$, $N_\mathrm{c,x}=3$,
$r_c\sim 2$), the sparse-ring cost has prefactor
$N_\mathrm{c,x}^{2}C_a^{3}\sim 2.4\times 10^{5}$; the factorised
kernel has prefactor $R\,r_c\,C_a + R^{2}\sim 60R + R^{2}$. For
$R=100$, the factorised kernel is roughly two orders of magnitude
cheaper.

\paragraph{Crossover.}
The factorised kernel is cheaper than the sparse ring when
\begin{equation}\label{eq:crossover}
  R\,r_c\,C_a + R^{2}
  \;<\; N_\mathrm{c,x}^{2}\,C_a^{3},
\end{equation}
which, for $R\ll C_a$, simplifies to $R\lesssim N_\mathrm{c,x}^{2}\,C_a^{2}/r_c\sim 10^{4}$
at the same parameters; for $R$ comparable to $C_a$, the $R^{2}$
term enters and the crossover moves to
$R\lesssim N_\mathrm{c,x}\,C_a^{3/2}\sim 500$. Reported CP and
Waring ranks for FC3 of Si, MgO, HgTe, TiNiSn, ZrO$_2$
\cite{luo_tensor_2025} are well below these bounds, placing the
factorised kernel firmly in the favourable regime.

\paragraph{Batching over $\tau$.}
The kernel is local in $\tau$: each $\tau$-slice produces an
independent $\Sigma(\tau)$. The natural batching strategy is to
process $N_b$ $\tau$-slices in a single call, stacking the
$N_b$ propagator slices into a $\ndof\times\ndof\times N_b$ tensor
and the corresponding $V$ matrices into a $\ndof\times R\times N_b$
tensor. The SpMM becomes a batched SpMM (cuSPARSE/MKL support);
the GEMM becomes a batched GEMM. Batch size $N_b$ is chosen
adaptively from the free HBM, with the dominant memory term being
the $V$ tensor at $\cO(N_b R\,\ndof)$.

\paragraph{Sparse-factor optimisation.}
When the factor vectors $W$ are intra-block sparse with density
$C_a/\NBS$ (the case for FC3-localised factors, see
\cref{ssec:phph_decomp_blocks}), the GEMM in
\cref{eq:kernel_step2} can be replaced by a sparse--dense product at
cost $\cO(R^{2}\,C_a\,N_\mathrm{c,x})$, reducing the second term in
\cref{eq:per_tau_cost} by a factor $\NBS/(C_a N_\mathrm{c,x})$. The
SpMM cost is unaffected (it is already proportional to
$\mathrm{nnz}(g)$, not to the density of $W$).

% ----------------------------------------------------------------------
\subsection{Transverse periodic case}
\label{ssec:phph_decomp_qspace}
% ----------------------------------------------------------------------

In the transverse-periodic setting of \cref{sub:transper_negf}, the
SSE \cref{eq:sigma_phph} carries an internal sum over a transverse
wavevector $\qp'$ which couples $g$ at $\qp'$ and $\qp-\qp'$ via the
FC3 evaluated at the same pair of momenta. We carry the
factorised kernel through the partial Fourier transform and identify
the structural simplifications.

\paragraph{Partial Fourier transform of CP factors.}
Performing the CP decomposition on the real-space FC3
$\Phi^{(3)}_{\mu,\,l_1\nu_1,\,l_2\nu_2}$ gives factor matrices in
which the second and third legs include the transverse cell label.
The partial Fourier transform \cref{eq:fc3_fourier} acts on these
cell labels and produces
\begin{equation}\label{eq:cp_fourier_decomp}
  \tilde\Phi^{(3)}_{\mu\nu_1\nu_2}(\qp,\qp')
  \;=\; \sum_{r}\lambda_r\,a_r(\mu)\,b_r(\nu_1;\qp)\,c_r(\nu_2;\qp'),
\end{equation}
with $\qp$-dependent factor vectors
$b_r(\nu_1;\qp)=\sum_{l_1}\tilde b_r(l_1,\nu_1)e^{-i\qp\cdot l_1}$ and
likewise for $c_r$. The CP rank is preserved exactly under the
partial Fourier transform: the decomposition
\cref{eq:fc3_separable_general} commutes with the cell-FT because the
FT acts on each factor leg separately. For Waring, the single
factor vector $u_r$ has all three legs Fourier-transformed against
the appropriate cell labels.

\paragraph{Factorised kernel in $\qp$-space.}
With $\Phi_L=\Phi_R$ at $(\qp',\qp-\qp')$ and $(\qp'-\qp,-\qp')$
respectively, and noting that for real factor vectors
$b_r(\nu;-\qp)=b_r(\nu;\qp)^{*}$, the kernel
\cref{eq:bubble_factorised} becomes
\begin{align}
  \label{eq:bubble_qspace}
  \mathcal{B}_{xx'}(\tau,\qp,\qp')
  &= \sum_{rs}\lambda_r\lambda_s\,a_r(x)\,a_s(x')\,
      P_{rs}(\tau,\qp,\qp')\,Q_{rs}(\tau,\qp,\qp'),\\
  \label{eq:P_qspace}
  P_{rs}(\tau,\qp,\qp')
  &= b_r(\cdot;\qp')^{\!\top}\,g(\tau,\qp')\,c_s(\cdot;\qp-\qp')^{*},\\
  \label{eq:Q_qspace}
  Q_{rs}(\tau,\qp,\qp')
  &= c_r(\cdot;\qp-\qp')^{\!\top}\,g(\tau,\qp-\qp')\,b_s(\cdot;\qp')^{*}.
\end{align}
The two bilinear forms now involve $g$ at \emph{different} transverse
momenta, $\qp'$ and $\qp-\qp'$, and are no longer related by transposition
even for Waring. They are, however, related by a change of summation
variable: under $\qp'\to\qp-\qp'$,
$P_{rs}(\tau,\qp,\qp-\qp')=Q_{rs}(\tau,\qp,\qp')$. After summing over
the full Monkhorst--Pack mesh (which is closed under
$\qp'\to\qp-\qp'$ for appropriate shifts), the contributions of $P$ and
$Q$ to $\Sigma$ are equal, so only $P$ need be evaluated.

\paragraph{Kernel restructure.}
The change-of-variable identity collapses
\cref{eq:bubble_qspace,eq:P_qspace,eq:Q_qspace} to
\begin{equation}\label{eq:sigma_qspace_final}
  \Sigma_{xx'}(\tau,\qp)
  = \frac{2}{N}\!\sum_{\qp'}\!
    \sum_{rs}\lambda_r\lambda_s a_r(x)a_s(x')\,
    P_{rs}(\tau,\qp,\qp')\,P_{rs}(\tau,\qp,\qp-\qp'),
\end{equation}
which is a convolution along $\qp'$ of the matrix-valued quantity
$P(\tau,\qp,\qp')$ with its shift $P(\tau,\qp,\qp-\qp')$. The
$\qp'$-sum is therefore a Fourier transform of a per-block product, and
the $N_q^{2}$ factor in \cref{eq:cost_q} is replaced by
$N_q \log N_q$ after a single FFT along the transverse axis. The
explicit factor of 2 in \cref{eq:sigma_qspace_final} comes from the
$P\leftrightarrow Q$ symmetry; for CP without internal symmetry the
factor is 1 and both bilinear forms enter explicitly.

% ----------------------------------------------------------------------
\subsection{Distributed implementation}
\label{ssec:phph_decomp_dist}
% ----------------------------------------------------------------------

The four-redistribution pipeline of \cref{ssec:phph_distribution} is
preserved bit-for-bit. The decomposition changes only Stage~2 of the
pipeline (the per-$\tau$ ring contraction). Stages~1 and~3
(FFT $G\to g$ and IFFT $\sigma\to\Sigma$), the two interconvertible
distributions \texttt{stack}/\texttt{nnz}, and the four
\texttt{stack}$\leftrightarrow$\texttt{nnz} dtransposes are unchanged.

\paragraph{Stage~2 with decomposition.}
The ring contraction
\cref{eq:phph_step_A,eq:phph_step_B,eq:phph_step_S} is replaced by
the unified kernel
\cref{eq:kernel_step1,eq:kernel_step2,eq:kernel_step3,eq:kernel_step4}.
On a single rank holding a $\tau$-slice of $g$ in the \texttt{stack}
distribution, the kernel is purely local in $\tau$ and consists of
SpMM, GEMM, and the sparse-output sandwich. No communication is
incurred inside Stage~2 except the FC3-dictionary halo discussed
below.

\paragraph{Per-rank arithmetic.}
Summing over $N_\tau$, $N_q^{2}$ (or $N_q\log N_q$ in the
$\qp'\to\qp-\qp'$ symmetrised form of
\cref{eq:sigma_qspace_final}) and distributing across $P=P_E\,P_S$
ranks gives
\begin{equation}\label{eq:cost_phph_dist_cp}
  W^{\mathrm{rank}}_\mathrm{phph,decomp}
  \;=\;
  \cO\!\Bigl(
    \tfrac{N_\tau\,N_q^{2}\,N_{AO}\,(R\,r_c\,C_a + R^{2})}{P}
  \Bigr).
\end{equation}
The $C_a^{3}\,N_\mathrm{c,x}^{2}$ factor of \cref{eq:cost_phph_dist}
is replaced by $R\,r_c\,C_a + R^{2}$. At $C_a=30$, $r_c=2$, $R=100$:
the old factor is $2.4\times 10^{5}$, the new factor is
$\sim 1.6\times 10^{4}$ --- a 15$\times$ arithmetic reduction. With
$\qp'$-FFT symmetrisation (\cref{ssec:phph_decomp_qspace}), an
additional factor $N_q/\log N_q\sim 50$ is gained for production
$N_q=400$.
 